%% =====================================================================
%%  Capítulo 5 -- Órbitas, osciladores y ciclos límite
%%  Formato nuevo (noviembre de 2026). Integra: Handout4a.tex (2024),
%%  LotkaVolterra.tex, el rediseño de teoría del índice y Nowak
%%  (septiembre de 2026), y las clases de 2026 de Manuela Chacón
%%  ("Ciclos": Duffing como motivación, órbitas cerradas, ciclos límite,
%%  dinámica de fase, osciladores acoplados, luciérnagas y Kuramoto,
%%  con la aplicación de sincronización resiliente).
%%  Correcciones respecto a las versiones anteriores: ver \S\ref{sec:oc-cierre}.
%%  Números verificados con herramientas/cap05_numeros.py (59/59) y
%%  codigo/cap05_orbitas_ciclos/verificarCap05.m (34/34 en Octave 8.4).
%% =====================================================================

\capitulo[cap:orbitas]{Órbitas, osciladores y ciclos límite}
  {Se escribe para atar al lector, para adueñarse de él, para seducirlo, para subyugarlo, para entrar en el espíritu de otro y quedarse allí [\dots]}
  {Enrique Vila-Matas}

\begin{objetivos}[partes/05-orbitas-ciclos.html]
  \item Distinguir una órbita cerrada de un ciclo límite y explicar por qué un sistema lineal no puede tener ciclos límite.
  \item Analizar flujos en el círculo y osciladores de fase: período, arrastre y deriva, y el umbral de sincronización del modelo de Kuramoto.
  \item Calcular el índice de una curva cerrada y usarlo, junto con los sistemas gradiente, las funciones de Lyapunov y el criterio de Bendixson--Dulac, para \emph{descartar} órbitas cerradas.
  \item Construir una región que atrape trayectorias y usar el teorema de Poincaré--Bendixson para \emph{probar} que existe una órbita periódica.
  \item Aplicar el mismo arsenal a modelos de poblaciones (Lotka--Volterra y sus variantes) y a un modelo de cambio de lenguaje.
\end{objetivos}

%% ---------------------------------------------------------------------
\section{¿Qué hay además de los equilibrios?}
\label{sec:oc-masalla}
%% ---------------------------------------------------------------------

En el \cref{cap:bidimensionales} estudiamos sistemas planos \emph{cerca} de sus equilibrios: linealización, Hartman--Grobman, variedad central. Es un análisis local: dice qué pasa en una vecindad de $\eq$ y nada sobre a dónde van las trayectorias que salen de ella. Después de los equilibrios, el comportamiento de largo plazo más simple es el \emph{periódico}: una trayectoria que vuelve al mismo estado, $x(t+T)=x(t)$. Sobre ellas hay dos preguntas, y todo el capítulo gira alrededor de las dos:
\begin{center}
\sffamily\large\color{snlazul}¿existen? \qquad ¿son estables?
\end{center}

\begin{ejemplo}[label=ej:oc-duffing]{Duffing sin forzamiento ni fricción}
Para $\ddot x+\delta\dot x+x^3-x=0$ con $\delta=0$: (a) escriba el modelo de estado; (b) halle los equilibrios; (c) clasifíquelos por linealización; (d) dibuje el plano de fase. ¿Explica la linealización lo que ve?
\solucion (a) Con $x_1=x$, $x_2=\dot x$: $\dot x_1=x_2$, $\dot x_2=x_1-x_1^3$. (b) $x_2=0$ y $x_1(1-x_1^2)=0$: equilibrios en $(-1,0)$, $(0,0)$ y $(1,0)$. (c) $\jac fx=\begin{bmatrix}0&1\\1-3x_1^2&0\end{bmatrix}$. En el origen los valores propios son $\pm1$: una silla (hiperbólica, Hartman--Grobman sí concluye). En $(\pm1,0)$ son $\pm j\sqrt2$: parte real nula, Hartman--Grobman \emph{no} concluye. (d) Multiplicando la ecuación por $\dot x$,
\[
  \frac{\dif}{\dif t}H=0,\qquad H(x_1,x_2)=\tfrac12x_2^2-\tfrac12x_1^2+\tfrac14x_1^4 .
\]
$H$ se conserva: las trayectorias viven en sus curvas de nivel. Alrededor de $(\pm1,0)$, donde $H$ tiene mínimos ($H=-\tfrac14$), las curvas de nivel son cerradas: son centros \emph{no lineales}, rodeados de una familia continua de órbitas cerradas. El nivel $H=0$ pasa por la silla y forma dos órbitas homoclínicas, que salen de la silla y vuelven a ella (\cref{fig:oc-duffing}).
\end{ejemplo}

\notamargen{Un sistema con $\dot x_1=\partial H/\partial x_2$, $\dot x_2=-\partial H/\partial x_1$ es \emph{hamiltoniano}; Duffing sin fricción lo es. Los sistemas hamiltonianos conservan $H$ y no pueden tener equilibrios asintóticamente estables ni ciclos límite.}
\script{cap05_orbitas_ciclos/C5_Ej1_Duffing.m}
\begin{figure}[ht]
\centering
\includegraphics[width=\linewidth]{cap05/f_duffing}
\caption{Izquierda: Duffing con $\delta=0$; las trayectorias son curvas de nivel de $H$. En rojo, las órbitas homoclínicas ($H=0$). Derecha: el período de las órbitas que rodean a $(1,0)$ depende de la órbita; solo cerca del centro se acerca a $2\pi/\sqrt2$, el que predice la linealización.}
\label{fig:oc-duffing}
\end{figure}

El ejemplo deja tres lecciones. La linealización no decide en los equilibrios no hiperbólicos; aquí decidió una \emph{primera integral} (una cantidad conservada). Las órbitas cerradas de Duffing no son ``una oscilación'': son infinitas, una por cada nivel de energía, y cada una con su propio período ($4.63$ para $H=-0.2$, $7.41$ para $H=-0.01$, y crece sin cota al acercarse a la homoclínica). Y basta un poco de fricción ($\delta>0$) para destruirlas todas (\cref{sec:oc-descartar}).

\codigorepo{cap05_orbitas_ciclos/C5_Ej1_Duffing.m}{energia}{$H$ se conserva a lo largo de la simulación}

\begin{trampa}[label=tr:oc-centro]{$\pm j\omega$ en la linealización no es un centro}
\plausible{En $(\pm1,0)$ los valores propios son $\pm j\sqrt2$; por lo tanto el equilibrio es un centro del sistema no lineal.}
En Duffing la conclusión es cierta, pero no por la linealización: es cierta porque $H$ se conserva. Sin una razón así, es falsa. El sistema
\[
  \dot x_1=-x_2+x_1(x_1^2+x_2^2),\qquad \dot x_2=x_1+x_2(x_1^2+x_2^2)
\]
tiene la misma linealización en el origen que un centro ($\pm j$), y sin embargo en polares $\dot r=r^3$, $\dot\theta=1$: el origen es un foco \emph{inestable} y las soluciones escapan en $t=1/(2r_0^2)$ (\cref{fig:oc-centrofalso}). Los términos de orden superior, que Hartman--Grobman permite ignorar cuando el equilibrio es hiperbólico, deciden todo cuando no lo es.
\correcto{Con valores propios imaginarios puros, la linealización no concluye nada. Hace falta un argumento no lineal: una cantidad conservada (centro), una función de Lyapunov (estable o inestable) o el cálculo en polares.}
\end{trampa}

\begin{figure}[ht]
\centering
\includegraphics[width=0.86\linewidth]{cap05/f_centro_falso}
\caption{La misma linealización (izquierda) y dos comportamientos no lineales distintos: Duffing tiene centros; el sistema de la \cref{tr:oc-centro} tiene un foco inestable (derecha).}
\label{fig:oc-centrofalso}
\end{figure}

\codigorepo{cap05_orbitas_ciclos/C5_Ej1_Duffing.m}{contraejemplo}{El centro falso: $\dot r=r^3$}

%% ---------------------------------------------------------------------
\section{Órbitas cerradas}
\label{sec:oc-orbitas}
%% ---------------------------------------------------------------------

En todo el capítulo, $\dot x=f(x)$ es invariante en el tiempo, $x\in\R^2$ (salvo que se diga otra cosa) y $f$ es localmente Lipschitz, de modo que por cada punto pasa una única solución (capítulo~3). Escribimos $\flujo t{x_0}$ para la solución que en $t=0$ está en $x_0$ (el \clave{flujo}).

\begin{definicion}[label=def:oc-orbita]{Solución periódica y órbita cerrada}
Una solución $x(t)$ es \clave{periódica} si existe $T>0$ con $x(t+T)=x(t)$ para todo $t$; el menor de esos $T$ es el \clave{período}. Si la solución periódica no es un equilibrio, su trayectoria $\gamma=\{x(t):0\le t\le T\}$ es una \clave{órbita cerrada}.
\end{definicion}

Por unicidad, una órbita cerrada no puede pasar por un equilibrio: si $x(t^*)=\eq$, la única solución que pasa por $\eq$ es la constante, y la trayectoria no sería periódica no constante. Por la misma razón no se puede cortar a sí misma: por el punto de cruce pasarían dos soluciones distintas. Una curva que se ve como un ``ocho'' alrededor de dos centros (\cref{fig:oc-duffing}, en rojo) no es una órbita cerrada: son dos órbitas homoclínicas más el equilibrio que las une, y cada una tarda un tiempo infinito en recorrerse.

\begin{proposicion}[label=prop:oc-lineal]{Las órbitas cerradas lineales vienen en familias}
Si $\dot x=Ax$ tiene una solución periódica no constante $x(t)$, entonces $c\,x(t)$ es solución periódica para todo $c\neq0$. Ninguna órbita cerrada de un sistema lineal está aislada.
\end{proposicion}

\begin{proof}
$\frac{\dif}{\dif t}(cx)=c\dot x=cAx=A(cx)$, y $cx(t+T)=cx(t)$. Las órbitas $c\gamma$ llenan una vecindad de $\gamma$.
\end{proof}

Esto tiene una consecuencia práctica: en un oscilador lineal (un LC ideal, un resorte sin fricción) la \emph{amplitud} de la oscilación la decide la condición inicial, y cualquier perturbación la cambia para siempre. Además, la oscilación es frágil: cualquier amortiguamiento, por pequeño que sea, la destruye, y cualquier término negativo la hace crecer sin cota. Un reloj, un corazón o un oscilador electrónico no pueden funcionar así.

%% ---------------------------------------------------------------------
\subsection{Conjuntos límite}
%% ---------------------------------------------------------------------

Para hablar de trayectorias que \emph{se acercan} a una órbita hace falta un lenguaje para el comportamiento de largo plazo.

\begin{definicion}[label=def:oc-omega]{Conjuntos $\omega$-límite y $\alpha$-límite}
Un punto $z$ es un \clave{punto $\omega$-límite} de $x_0$ si existe una sucesión $t_k\to\infty$ con $\flujo{t_k}{x_0}\to z$. El conjunto de todos ellos es $\omegalim{x_0}$. Con $t_k\to-\infty$ se definen los puntos y el conjunto \clave{$\alpha$-límite}, $\alphalim{x_0}$.
\end{definicion}

\notamargen{$\omega$ y $\alpha$ son la última y la primera letra del alfabeto griego: el final y el comienzo de la trayectoria.}
Si $x_0$ es un equilibrio, $\omegalim{x_0}=\{x_0\}$. Si $x_0$ está en una órbita cerrada $\gamma$, $\omegalim{x_0}=\gamma$. Si la solución converge a un equilibrio $\eq$, $\omegalim{x_0}=\{\eq\}$. El caso interesante es el de una trayectoria que \emph{no está} en $\gamma$ y tiene $\omegalim{x_0}=\gamma$: se enrolla sobre la órbita sin alcanzarla nunca.

\begin{lema}[label=lem:oc-omega]{Propiedades de $\omegalim{x_0}$ \citep[Lema~4.1]{Kha02}}
Si la solución $\flujo t{x_0}$ es acotada para $t\ge0$, entonces $\omegalim{x_0}$ es no vacío, compacto, invariante (si $z\in\omegalim{x_0}$, toda la trayectoria de $z$ está en $\omegalim{x_0}$) y conexo, y $\flujo t{x_0}$ se acerca a $\omegalim{x_0}$ cuando $t\to\infty$.
\end{lema}

Este lema es la pieza que hace funcionar el teorema de Poincaré--Bendixson (\cref{sec:oc-pb}) y, en el capítulo~8, el principio de invarianza de LaSalle.

%% ---------------------------------------------------------------------
\section{Ciclos límite}
\label{sec:oc-ciclos}
%% ---------------------------------------------------------------------

\begin{definicion}[label=def:oc-ciclo]{Ciclo límite}
Un \clave{ciclo límite} es una órbita cerrada \emph{aislada}: alguna vecindad suya no contiene otras órbitas cerradas. Equivalentemente (en el plano), es una órbita cerrada $\gamma$ que es el conjunto $\omega$-límite o $\alpha$-límite de algún punto que no está en $\gamma$.
\end{definicion}

Todo ciclo límite es una órbita cerrada; no toda órbita cerrada es un ciclo límite (las de Duffing y las de un sistema lineal no lo son). Por la \cref{prop:oc-lineal}, los ciclos límite son un fenómeno exclusivamente no lineal. Las trayectorias vecinas a un ciclo límite se acercan o se alejan de él en espiral, y eso da tres tipos (\cref{fig:oc-tipos}): \clave{estable} (se acercan por los dos lados), \clave{inestable} (se alejan por los dos lados) y \clave{semiestable} (se acercan por un lado y se alejan por el otro).

\begin{figure}[ht]
\centering
\includegraphics[width=\linewidth]{cap05/f_tipos_ciclo}
\caption{Ciclos límite estable ($\dot r=r(1-r^2)$), inestable ($\dot r=-r(1-r^2)$) y semiestable ($\dot r=r(1-r)^2$), todos con $\dot\theta=1$. En rojo, trayectorias que empiezan adentro; en azul, afuera.}
\label{fig:oc-tipos}
\end{figure}

\begin{ejemplo}[label=ej:oc-polar]{Un ciclo límite que se puede calcular}
Analice $\dot x_1=-x_2+x_1(1-x_1^2-x_2^2)$, $\dot x_2=x_1+x_2(1-x_1^2-x_2^2)$.
\solucion El único equilibrio es el origen, con jacobiano $\begin{bmatrix}1&-1\\1&1\end{bmatrix}$ y valores propios $1\pm j$: un foco inestable. La linealización no dice a dónde van las trayectorias. En polares, $r^2=x_1^2+x_2^2$ y $\theta=\operatorname{atan2}(x_2,x_1)$, se usan las identidades
\begin{equation}
  r\dot r=x_1\dot x_1+x_2\dot x_2,\qquad r^2\dot\theta=x_1\dot x_2-x_2\dot x_1,
  \label{eq:oc-polares}
\end{equation}
válidas en todo punto con $r>0$ (no hace falta excluir $x_1=0$). Aquí dan $\dot r=r(1-r^2)$ y $\dot\theta=1$: el radio y el ángulo se desacoplan. La ecuación radial es la logística del capítulo~2, con $r=0$ inestable y $r=1$ estable, y tiene solución explícita
\[
  r(t)=\Bigl(1+\bigl(r_0^{-2}-1\bigr)e^{-2t}\Bigr)^{-1/2}.
\]
Toda trayectoria con $r_0>0$ se enrolla sobre el círculo unitario, que es un ciclo límite estable de período $2\pi$ (\cref{fig:oc-linealciclo}, derecha). Desde $r_0=0.1$, $r(5)=0.99776$; la simulación coincide con la fórmula con error menor que $10^{-10}$.
\end{ejemplo}

\codigorepo{cap05_orbitas_ciclos/C5_Ej2_CicloPolar.m}{ciclo}{La solución en polares contra la simulación}

\begin{figure}[ht]
\centering
\includegraphics[width=0.86\linewidth]{cap05/f_lineal_vs_ciclo}
\caption{Izquierda: el oscilador armónico tiene una familia continua de órbitas; la amplitud la escoge la condición inicial. Derecha: el ciclo límite del \cref{ej:oc-polar} atrae a las trayectorias de adentro y de afuera; la amplitud la escoge el sistema.}
\label{fig:oc-linealciclo}
\end{figure}

%% ---------------------------------------------------------------------
\subsection{Oscilaciones autosostenidas: de Rayleigh a Liénard}
%% ---------------------------------------------------------------------

Un ciclo límite estable es una \clave{oscilación autosostenida}: el sistema oscila sin ninguna entrada periódica, con un período, una amplitud y una forma de onda propios, y si se le perturba, vuelve a ellos. El latido del corazón, el ritmo circadiano de la temperatura corporal, las reacciones químicas que oscilan espontáneamente (Belousov--Zhabotinsky) y el aleteo (\emph{flutter}) de un ala son ejemplos.

Los primeros modelos vinieron de la acústica y de la electrónica. Lord Rayleigh, estudiando las vibraciones \emph{mantenidas} ---la cuerda de un violín frotada por el arco, el tubo de un órgano---, propuso en 1883 la ecuación $\ddot x+\epsilon\bigl(\tfrac13\dot x^2-1\bigr)\dot x+x=0$, es decir,
\[
  \dot x_1=x_2,\qquad \dot x_2=-x_1-\epsilon\bigl(\tfrac13x_2^2-1\bigr)x_2 \qquad\citep{Ray1883}.
\]
Con los tubos de vacío aparecieron oscilaciones de amplitud fija que no dependían del arranque, y Balthasar van der Pol propuso en 1926 la ecuación que hoy lleva su nombre \citep{van26}:
\begin{equation}
  \ddot x+\mu(x^2-1)\dot x+x=0,\qquad\mu>0 .
  \label{eq:oc-vdp}
\end{equation}
Es un oscilador armónico con un amortiguamiento que \emph{cambia de signo}: para $\abs x<1$ el coeficiente $\mu(x^2-1)$ es negativo y el sistema inyecta energía (las oscilaciones pequeñas crecen); para $\abs x>1$ es positivo y la disipa (las grandes decrecen). Entre las dos, una oscilación de amplitud cercana a $2$ (\cref{fig:oc-vdp}). Para $\mu$ grande la oscilación es de \emph{relajación}: lentas cargas y descargas rápidas, con período aproximadamente $(3-2\ln2)\mu$.

\video{V5_3_vanderpol}{Todas las trayectorias llegan al mismo ciclo}
\begin{figure}[ht]
\centering
\includegraphics[width=\linewidth]{cap05/f_vanderpol}
\caption{Van der Pol. Con $\mu=1$ el ciclo tiene período $6.663$ y amplitud $2.009$; con $\mu=5$, período $11.61$ y la forma de onda de una oscilación de relajación.}
\label{fig:oc-vdp}
\end{figure}

¿Cómo se \emph{prueba} que \eqref{eq:oc-vdp} tiene un ciclo límite, y que es uno solo? Alfred Liénard lo hizo en 1928 para una familia más amplia \citep{Lie28}.

\begin{teorema}[label=thm:oc-lienard]{Teorema de Liénard \citep[Teorema~7.4.1]{Str15}}
Considere $\ddot x+f(x)\dot x+g(x)=0$, es decir, $\dot x_1=x_2$, $\dot x_2=-g(x_1)-f(x_1)x_2$, y suponga que
\begin{enumerate}
  \item $f$ y $g$ son $\Cont1$;
  \item $g$ es impar y $g(x)>0$ para $x>0$;
  \item $f$ es par;
  \item $F(x)=\int_0^xf(u)\,\dif u$ tiene exactamente un cero positivo en $x=a$, es negativa en $(0,a)$, positiva y no decreciente para $x>a$, y $F(x)\to\infty$ cuando $x\to\infty$.
\end{enumerate}
Entonces el sistema tiene un único ciclo límite, que rodea el origen y es estable.
\end{teorema}

Para van der Pol, $f(x)=\mu(x^2-1)$ es par, $g(x)=x$ es impar y positiva para $x>0$, y $F(x)=\mu(\tfrac13x^3-x)=\tfrac\mu3x(x^2-3)$ tiene su único cero positivo en $a=\sqrt3$, es negativa en $(0,\sqrt3)$ y creciente después. El teorema da existencia, unicidad y estabilidad del ciclo para todo $\mu>0$. Es, en este capítulo, el único resultado que dice \emph{cuántos} ciclos hay.

\codigorepo{cap05_orbitas_ciclos/C5_Ej3_VanDerPol.m}{ciclo}{Período y amplitud del ciclo de van der Pol}

\begin{trampa}[label=tr:oc-orbital]{Un ciclo estable no es una solución asintóticamente estable}
\plausible{Las trayectorias vecinas convergen al ciclo límite; por lo tanto la solución periódica es asintóticamente estable en el sentido de Lyapunov: si $x(0)$ está cerca de $\bar x(0)$, entonces $x(t)-\bar x(t)\to0$.}
Lo que converge es la \emph{órbita}, no el instante en que se pasa por cada punto. Una trayectoria que empieza un poco afuera del ciclo de van der Pol llega al ciclo con un desfase: $x(t)$ se acerca al conjunto $\gamma$, pero $x(t)-\bar x(t)$ no tiende a cero, salvo que se corrija la fase. Basta comparar dos puntos \emph{sobre} el mismo ciclo: siguen separados para siempre, aunque los dos estén en el ciclo.
\correcto{Un ciclo límite estable es \emph{orbitalmente} asintóticamente estable: $\dist(x(t),\gamma)\to0$. Cada trayectoria vecina converge, además, a la solución periódica con alguna \emph{fase asintótica} propia. La estabilidad de Lyapunov de una solución periódica de un sistema invariante en el tiempo nunca es asintótica, porque un corrimiento de fase sobre el ciclo no decae.}
\end{trampa}

%% ---------------------------------------------------------------------
\section{Flujos en el círculo y osciladores de fase}
\label{sec:oc-circulo}
%% ---------------------------------------------------------------------

El modelo más simple de algo que oscila no necesita dos dimensiones. Basta una variable angular,
\begin{equation}
  \dot\theta=f(\theta),\qquad\theta\in S^1,
  \label{eq:oc-circulo}
\end{equation}
donde $S^1$ es el círculo: $\theta$ y $\theta+2\pi$ son el \emph{mismo} punto. La dinámica sigue siendo unidimensional, pero ahora una trayectoria puede dar la vuelta y volver a su punto de partida: hay soluciones periódicas sin equilibrios, algo imposible en la recta (\cref{prop:un-monotonia}). Para que \eqref{eq:oc-circulo} defina un campo en el círculo, $f$ tiene que asignar una sola velocidad a cada punto, es decir, ser $2\pi$-periódica: $f(\theta+2\pi)=f(\theta)$. $\dot\theta=\theta$ \emph{no} es un campo en el círculo, porque daría velocidades distintas en $\theta=0$ y en $\theta=2\pi$.

\begin{ejemplo}[label=ej:oc-seno]{El campo $\dot\theta=\sin\theta$}
Dibuje el flujo en el círculo (con $\theta$ creciendo en sentido antihorario) y clasifique los equilibrios.
\solucion Equilibrios: $\sin\theta^*=0$, es decir, $\theta^*=0$ y $\theta^*=\pi$. En $(0,\pi)$, $\dot\theta>0$: el flujo va en sentido antihorario; en $(\pi,2\pi)$, en sentido horario. Las flechas salen de $0$ y llegan a $\pi$: $\theta^*=0$ es inestable y $\theta^*=\pi$ estable. Con la derivada: $f'(0)=\cos0=1>0$ y $f'(\pi)=-1<0$ (\cref{fig:oc-circulo}).
\end{ejemplo}

\begin{figure}[ht]
\centering
\includegraphics[width=0.92\linewidth]{cap05/f_circulo_seno}
\caption{$\dot\theta=\sin\theta$ dibujado en la recta (identificando $0$ con $2\pi$) y en el círculo. En el círculo, los equilibrios estables e inestables tienen que alternarse.}
\label{fig:oc-circulo}
\end{figure}

%% ---------------------------------------------------------------------
\subsection{El oscilador uniforme}
%% ---------------------------------------------------------------------

Un punto del círculo es una \clave{fase}. El oscilador más simple es aquel cuya fase avanza a velocidad constante: $\dot\theta=\omega$, con solución $\theta(t)=\omega t+\theta_0$ y período $T=2\pi/\omega$. No hay amplitud; para eso hace falta una segunda dimensión. Lo que sí hay es la posibilidad de comparar dos osciladores por su \emph{diferencia de fase}.

\begin{ejemplo}[label=ej:oc-corredores]{Dos corredores en una pista}
$P_1$ da una vuelta a la pista en $T_1=4$ minutos y $P_2$ en $T_2=5$. Si salen juntos, ¿cuánto tarda $P_1$ en sacarle una vuelta a $P_2$?
\solucion Con $\dot\theta_i=\omega_i=2\pi/T_i$, la diferencia de fase $\varphi=\theta_1-\theta_2$ cumple $\dot\varphi=\omega_1-\omega_2$. $P_1$ saca una vuelta cuando $\varphi$ llega a $2\pi$:
\[
  T_{\text{vuelta}}=\frac{2\pi}{\omega_1-\omega_2}=\Bigl(\frac1{T_1}-\frac1{T_2}\Bigr)^{-1}=\Bigl(\frac14-\frac15\Bigr)^{-1}=20\ \text{minutos}.
\]
Es el fenómeno de \emph{batido}: dos osciladores de frecuencias cercanas se ponen en fase y se desfasan periódicamente, con una frecuencia igual a la diferencia de las suyas. Así suenan dos campanas o dos cuerdas casi afinadas.
\end{ejemplo}

\notamargen{\textbf{Notación.} En este capítulo $\varphi$ es una diferencia de fase y $\phi(t;x_0)$ el flujo. No se confundan.}

%% ---------------------------------------------------------------------
\subsection{El oscilador no uniforme}
%% ---------------------------------------------------------------------

La ecuación
\begin{equation}
  \dot\theta=\omega-a\sin\theta,\qquad\omega>0,\ a\ge0,
  \label{eq:oc-nouniforme}
\end{equation}
aparece en los lazos de enganche de fase (PLL) en electrónica, en el disparo periódico de neuronas, en el ciclo sueño--vigilia y, como veremos, en las luciérnagas. Con $a=0$ es el oscilador uniforme; $a$ introduce una no uniformidad: el oscilador va rápido cerca de $\theta=-\pi/2$ y lento cerca de $\theta=\pi/2$. Hay tres casos (\cref{fig:oc-nouniforme}):
\begin{itemize}
  \item $a<\omega$: $\dot\theta>0$ siempre, no hay equilibrios y la fase da vueltas, cada vez más despacio al cruzar $\pi/2$. El período es
  \begin{equation}
    T=\int_0^{2\pi}\frac{\dif\theta}{\omega-a\sin\theta}=\frac{2\pi}{\sqrt{\omega^2-a^2}} .
    \label{eq:oc-periodo}
  \end{equation}
  \item $a=\omega$: aparece un equilibrio semiestable en $\theta=\pi/2$ (atrae por la izquierda y repele por la derecha).
  \item $a>\omega$: dos equilibrios, $\sin\theta^*=\omega/a$; el de $\cos\theta^*>0$ es estable y el otro inestable. La fase se \emph{engancha}.
\end{itemize}
Cuando $a\to\omega^-$ el período \eqref{eq:oc-periodo} crece sin cota, como $1/\sqrt{\omega-a}$: con $\omega=2$, vale $3.63$ para $a=1$, $10.06$ para $a=1.9$ y $31.46$ para $a=1.99$. Ese ``cuello de botella'' es la huella de una bifurcación silla-nodo en el círculo, que veremos en el capítulo~6.

\begin{figure}[ht]
\centering
\includegraphics[width=\linewidth]{cap05/f_no_uniforme}
\caption{$\dot\theta=\omega-a\sin\theta$ con $\omega=2$: deriva sin equilibrios ($a=1$), un equilibrio semiestable ($a=2$) y un par estable--inestable ($a=3$). A la derecha, el período \eqref{eq:oc-periodo}, que explota al acercarse $a$ a $\omega$.}
\label{fig:oc-nouniforme}
\end{figure}

\codigorepo{cap05_orbitas_ciclos/C5_Ej4_Circulo.m}{nouniforme}{El período del oscilador no uniforme, simulado y con la fórmula}

\begin{pausa}
En el círculo, ¿puede un campo continuo tener exactamente un equilibrio estable y ningún otro equilibrio? ¿Y exactamente un equilibrio, semiestable?
\respuesta{Un solo equilibrio estable, no: si $f$ cambia de signo de $+$ a $-$ en $\theta^*$, al recorrer el círculo tiene que volver a cambiar de $-$ a $+$, y ese otro cero es un equilibrio inestable. En el círculo los cambios de signo vienen en pares. Un solo equilibrio semiestable, sí: $\dot\theta=1-\sin\theta$ tiene solo $\theta^*=\pi/2$, donde $f$ toca el cero sin cambiar de signo. Es el caso $a=\omega$.}
\end{pausa}

%% ---------------------------------------------------------------------
\section{Osciladores acoplados: arrastre y deriva}
\label{sec:oc-acoplados}
%% ---------------------------------------------------------------------

Hasta ahora cada oscilador evoluciona solo. En muchos sistemas los osciladores \emph{interactúan}, y la pregunta es si la interacción los sincroniza.

\begin{historia}[De los relojes de Huygens a Kuramoto]
En febrero de 1665, enfermo en cama, Christiaan Huygens notó que dos relojes de péndulo colgados de la misma viga terminaban oscilando en antifase, y que volvían a ella si se les perturbaba; lo llamó la ``simpatía de los relojes'' y lo publicó en el \emph{Horologium Oscillatorium} de 1673 \citep{Huy73}. Rayleigh describió en \emph{The Theory of Sound} la sincronización de tubos de órgano vecinos. En 1920, Eccles y Vincent patentaron el acoplamiento de dos generadores de triodo con frecuencias ligeramente distintas, que terminaban oscilando a la misma; Appleton y van der Pol extendieron esos experimentos y empezaron su teoría. En biología, de Mairan observó en 1729 que las hojas de una planta seguían abriéndose y cerrándose con el ritmo del día aun en la oscuridad: un reloj interno. Wiener se preguntó en los años cincuenta cómo se sincronizan las neuronas que producen el ritmo alfa del cerebro \citep{Wie58}; Winfree propuso en 1967 el primer modelo matemático de una población de osciladores biológicos acoplados y descubrió una transición de la incoherencia a la sincronía \citep{Win67}; y Kuramoto, en 1975, simplificó ese modelo hasta hacerlo resoluble \citep{Kur75,Kur84}. La revisión de \citet{DoBu14} cuenta la historia completa desde el punto de vista de control.
\end{historia}

Una población de $N$ osciladores de fase acoplados se modela, en general, como
\begin{equation}
  \dot\theta_i=\omega_i-\sum_{j=1}^Na_{ij}\sin(\theta_i-\theta_j),\qquad i=1,\dots,N,
  \label{eq:oc-redsin}
\end{equation}
donde $\omega_i$ es la frecuencia natural del oscilador $i$ y $a_{ij}\ge0$ el peso de la arista $j\to i$ de un grafo. El término $-a_{ij}\sin(\theta_i-\theta_j)$ frena al oscilador $i$ si va adelante de $j$ y lo acelera si va atrás. Linealizado alrededor de $\theta_i=\theta_j$, \eqref{eq:oc-redsin} es el protocolo de consenso que se usa en control de sistemas multiagente.

Con dos osciladores, la diferencia de fase $\varphi=\theta_1-\theta_2$ cumple
\begin{equation}
  \dot\varphi=(\omega_1-\omega_2)-(a_{12}+a_{21})\sin\varphi :
  \label{eq:oc-dos}
\end{equation}
toda la pregunta de sincronización se reduce al oscilador no uniforme \eqref{eq:oc-nouniforme}. Hay enganche si y solo si el acoplamiento total $a_{12}+a_{21}$ supera la diferencia de frecuencias $\abs{\omega_1-\omega_2}$.

%% ---------------------------------------------------------------------
\subsection{Luciérnagas}
%% ---------------------------------------------------------------------

\begin{quote}\small\itshape
Some twenty years ago I saw, or thought I saw, a synchronal or simultaneous flashing of fireflies. I could hardly believe my eyes, for such a thing to occur among insects is certainly contrary to all natural laws.
\par\hfill\normalfont--- Philip Laurent, carta a \emph{Science}, 1917 \citep{Lau17}
\end{quote}

En algunos lugares del sudeste asiático, miles de luciérnagas macho se posan en los árboles al atardecer y, poco a poco, empiezan a destellar al unísono. Llegan con fases distintas; nadie dirige; cada una ve los destellos de las vecinas y ajusta su ritmo. El experimento que permitió modelarlo fue sencillo: encender una linterna periódicamente frente a una luciérnaga \citep{Han78}. Si la frecuencia de la linterna es cercana a la natural de la luciérnaga (período cercano a $0.9$\,s), la luciérnaga la sigue: queda \clave{arrastrada} (\emph{entrained}). Si es demasiado rápida o demasiado lenta, la luciérnaga no la alcanza y la diferencia de fase avanza sin cesar: \clave{deriva de fase} (\emph{phase drift}).

\begin{ejemplo}[label=ej:oc-luciernaga]{Arrastre y deriva}
Sea $\theta$ la fase de la luciérnaga, con frecuencia natural $\omega$, y $\psi$ la de la linterna, $\dot\psi=\Omega$. Suponga que la luciérnaga acelera si el estímulo va adelante y frena si va atrás \citep{ErRi84}:
\[
  \dot\theta=\omega+A\sin(\psi-\theta),\qquad A>0 .
\]
Halle el rango de frecuencias $\Omega$ que arrastran a la luciérnaga, la diferencia de fase en el arrastre y el período de la deriva fuera de él.
\solucion La diferencia $\varphi=\psi-\theta$ cumple
\[
  \dot\varphi=\Omega-\omega-A\sin\varphi,
\]
otra vez \eqref{eq:oc-nouniforme}. Hay equilibrios si y solo si $\abs{\Omega-\omega}\le A$: el \clave{rango de arrastre} es $\omega-A\le\Omega\le\omega+A$. Dentro de él, el equilibrio estable es $\varphi^*=\arcsin\bigl((\Omega-\omega)/A\bigr)\in[-\pi/2,\pi/2]$ (allí $-A\cos\varphi^*<0$). Fuera de él, la deriva tiene período $T_{\text{deriva}}=2\pi/\sqrt{(\Omega-\omega)^2-A^2}$. Con $\omega=1$ y $A=0.5$: $\Omega=1.3$ arrastra con $\varphi^*=0.644$; $\Omega=1.6$ deriva con período $18.94$; $\Omega=2$, con período $7.26$ (\cref{fig:oc-luciernagas}). Midiendo el rango de arrastre en el laboratorio se estima $A$, la ``fuerza de reajuste'' de la luciérnaga.
\end{ejemplo}

\script{cap05_orbitas_ciclos/C5_Ej5_Luciernagas.m}
\begin{figure}[ht]
\centering
\includegraphics[width=\linewidth]{cap05/f_luciernagas}
\caption{Izquierda: dentro del rango de arrastre la diferencia de fase se queda quieta (verde); fuera, avanza a saltos, lenta cerca de $\varphi=\pi/2$. Derecha: el período de la deriva explota en los bordes del rango de arrastre.}
\label{fig:oc-luciernagas}
\end{figure}

\codigorepo{cap05_orbitas_ciclos/C5_Ej5_Luciernagas.m}{arrastre}{Arrastre y deriva de una luciérnaga}

\begin{trampa}[label=tr:oc-enfase]{Arrastrada no quiere decir en fase}
\plausible{Si la luciérnaga queda arrastrada por la linterna, destellan al mismo tiempo.}
En el arrastre la diferencia de fase es \emph{constante}, no nula: $\varphi^*=\arcsin\bigl((\Omega-\omega)/A\bigr)$, que es cero solo si $\Omega=\omega$. Con $\Omega=1.3$, $\omega=1$ y $A=0.5$, la luciérnaga destella siempre $0.64$\,rad (una décima de ciclo) después de la linterna. Tienen la misma \emph{frecuencia}, no la misma fase.
\correcto{Distinga \emph{sincronización de frecuencia} (las fases avanzan al mismo ritmo, $\dot\varphi\to0$) de \emph{sincronización de fase} ($\varphi\to0$). La primera exige que el acoplamiento compense la diferencia de frecuencias; la segunda, además, frecuencias iguales.}
\end{trampa}

\notamargen{La sesión de 2026 sobre osciladores se dictó el 31 de agosto, aniversario de la muerte de Jan C.\ Willems (1939--2013). Su teoría de los sistemas disipativos es la base del capítulo~11, donde volvemos a interconectar osciladores, esta vez contando la energía que intercambian.}

%% ---------------------------------------------------------------------
\section{El modelo de Kuramoto}
\label{sec:oc-kuramoto}
%% ---------------------------------------------------------------------

\notamargen{Esta sección es fenomenológica: se describe lo que pasa y por qué es plausible. Las pruebas de sincronización usan herramientas de Lyapunov y de pasividad (capítulos~7 a~11).}
Kuramoto tomó \eqref{eq:oc-redsin} con todos contra todos y el mismo peso:
\begin{equation}
  \dot\theta_i=\omega_i+\frac KN\sum_{j=1}^N\sin(\theta_j-\theta_i),\qquad i=1,\dots,N,
  \label{eq:oc-kuramoto}
\end{equation}
donde $K\ge0$ es la intensidad del acoplamiento y las frecuencias naturales $\omega_i$ se sacan de una distribución con densidad $g(\omega)$, simétrica alrededor de su centro (que, cambiando de marco de referencia, se puede tomar como cero). La pregunta: ¿a partir de qué $K$ emerge la sincronía?

Para medirla se usa el \clave{parámetro de orden}
\begin{equation}
  r\,e^{j\Psi}=\frac1N\sum_{j=1}^Ne^{j\theta_j} .
  \label{eq:oc-orden}
\end{equation}
Es el centroide de las $N$ fases vistas como puntos del círculo unitario: $r\in[0,1]$ mide la coherencia ($r\approx0$, fases repartidas; $r=1$, todas iguales) y $\Psi$ es la fase promedio. Multiplicando \eqref{eq:oc-orden} por $e^{-j\theta_i}$ y tomando la parte imaginaria, $r\sin(\Psi-\theta_i)=\frac1N\sum_j\sin(\theta_j-\theta_i)$, y el modelo queda
\begin{equation}
  \dot\theta_i=\omega_i+K\,r\sin(\Psi-\theta_i) .
  \label{eq:oc-campomedio}
\end{equation}
Cada oscilador ya no ve a los demás uno por uno: ve un solo ``estímulo'' de fase $\Psi$ e intensidad $Kr$. Es la luciérnaga frente a la linterna, \eqref{eq:oc-dos}, con la diferencia de que la linterna la fabrican entre todos. El mecanismo es una retroalimentación: un poco de coherencia ($r>0$) hace que más osciladores se enganchen (los que cumplen $\abs{\omega_i}\le Kr$), y eso aumenta $r$.

Kuramoto mostró que, cuando $N\to\infty$, la solución incoherente $r=0$ pierde la estabilidad en
\begin{equation}
  K_c=\frac2{\pi g(0)},
  \label{eq:oc-Kc}
\end{equation}
y que, para la distribución de Lorentz $g(\omega)=\frac{\gamma}{\pi(\gamma^2+\omega^2)}$, la coherencia de régimen es exactamente $r=\sqrt{1-K_c/K}$ para $K>K_c=2\gamma$ \citep{Kur84}. La \cref{fig:oc-kuramoto} compara esa predicción con una simulación de $N=400$ osciladores con $\gamma=0.5$: $r=0.58$, $0.71$ y $0.82$ para $K=1.5$, $2$ y $3$, contra $0.577$, $0.707$ y $0.816$ de la teoría.

\video{V5_1_kuramoto}{Las fases se agrupan; algunas siguen dando vueltas}
\script{cap05_orbitas_ciclos/C5_Ej6_Kuramoto.m}
\begin{figure}[ht]
\centering
\includegraphics[width=\linewidth]{cap05/f_kuramoto}
\caption{Kuramoto con $N=400$ y frecuencias de Lorentz ($\gamma=0.5$, $K_c=1$). Las fases en $t=100$ (color: frecuencia natural) para $K=0.5$ y $K=2$, con la flecha $re^{j\Psi}$. Derecha: $r$ de régimen contra la predicción $\sqrt{1-K_c/K}$.}
\label{fig:oc-kuramoto}
\end{figure}

\codigorepo{cap05_orbitas_ciclos/C5_Ej6_Kuramoto.m}{kuramoto}{Parámetro de orden de régimen para varios $K$}

\begin{trampa}[label=tr:oc-kuramoto]{Acoplar no es igualar}
\plausible{Con $K$ suficientemente grande, todas las fases convergen a la misma: $r\to1$.}
Con frecuencias naturales distintas, las fases \emph{no} pueden ser iguales: si $\theta_i=\theta_j$ para todo $i,j$, entonces $\dot\theta_i=\omega_i$ y se separan de inmediato. Lo que se alcanza es un estado en que los osciladores enganchados giran a una frecuencia común con desfases fijos, y para Lorentz quedan además osciladores rápidos sin enganchar (los de $\abs{\omega_i}>Kr$), que siguen dando vueltas. Aun con $K=3$, $r=0.82$.
\correcto{Para $K>K_c$, $r$ crece con $K$ pero se queda por debajo de $1$; solo $r\to1$ cuando $K\to\infty$. Igual que en la \cref{tr:oc-enfase}: el acoplamiento da sincronía de frecuencia, no de fase, salvo que las frecuencias sean idénticas.}
\end{trampa}

\begin{pausa}
Con dos osciladores ($N=2$) y frecuencias $\pm\Delta/2$, ¿qué valor de $K$ en \eqref{eq:oc-kuramoto} los engancha?
\respuesta{Con $N=2$, $a_{12}=a_{21}=K/2$ en \eqref{eq:oc-redsin}, y \eqref{eq:oc-dos} da $\dot\varphi=\Delta-K\sin\varphi$. Hay enganche si y solo si $K\ge\Delta$.}
\end{pausa}

\begin{profundiza}[Investigación: sincronización resiliente con pocos actuadores]
En una red real (una red eléctrica, una red de sensores) hay perturbaciones y se puede actuar sobre algunos nodos:
\[
  \dot\theta_i=\omega_i+\frac KN\sum_{j=1}^Na_{ij}\sin(\theta_j-\theta_i)+u_i(t)+d_i(t),
\]
con $d_i$ una disrupción y $u_i$ un control. ¿Cómo recuperar la sincronía después de una disrupción controlando solo unos pocos nodos y gastando poca energía? Es la pregunta del proyecto de Chacón-Chamorro, Charria-Macías, Giraldo y Quijano sobre sincronización resiliente de redes de Kuramoto \citep{CCGQ26}: una arquitectura en la que una red neuronal escoge en línea qué nodos actuar (\emph{pinning}) y un controlador de consenso los lleva de vuelta al estado sincronizado.
\end{profundiza}

%% ---------------------------------------------------------------------
\section{Teoría del índice}
\label{sec:oc-indice}
%% ---------------------------------------------------------------------

En Duffing (\cref{fig:oc-duffing}) las órbitas cerradas rodean a los centros, nunca a la silla sola; y la curva grande que rodea los tres equilibrios también es una órbita. ¿Es casualidad? ¿Debe una órbita cerrada rodear algún equilibrio? ¿Cuáles? La linealización no puede responder: mira un equilibrio a la vez. Necesitamos un número que se lea \emph{sobre} una curva y diga algo de lo que hay \emph{adentro}, sin mirar adentro.

%% ---------------------------------------------------------------------
\subsection{El número de vueltas del campo}
%% ---------------------------------------------------------------------

Sea $\Gamma$ una curva cerrada simple (no se corta a sí misma), recorrida en sentido antihorario, que \emph{no pasa por ningún equilibrio}. En cada punto $x\in\Gamma$ el campo $f(x)\neq0$ forma un ángulo $\theta_f(x)$ con el eje $x_1$. Al dar la vuelta a $\Gamma$, el vector $f$ gira continuamente y, al volver al punto de partida, vuelve a su dirección inicial: su ángulo cambió en un múltiplo entero de $2\pi$.

\begin{definicion}[label=def:oc-indice]{Índice de una curva}
El \clave{índice} de $\Gamma$ respecto a $f$ es el número neto de vueltas antihorarias que da el vector $f$ al recorrer $\Gamma$ una vez en sentido antihorario:
\[
  I_f(\Gamma)=\frac1{2\pi}\bigl[\theta_f\bigr]_\Gamma=\frac1{2\pi}\oint_\Gamma\frac{f_1\,\dif f_2-f_2\,\dif f_1}{f_1^2+f_2^2} .
\]
\end{definicion}

\notamargen{La integral es $\oint\dif\theta_f$ con $\theta_f=\arctan(f_2/f_1)$ derivado sin preocuparse por las ramas del arcotangente: $\dif\theta_f=(f_1\dif f_2-f_2\dif f_1)/(f_1^2+f_2^2)$.}
Para calcularlo a mano basta numerar unos cuantos puntos de $\Gamma$, dibujar en cada uno el vector $f$, \emph{trasladar} todos esos vectores al origen sin rotarlos, y contar cuántas vueltas dan en orden (\cref{fig:oc-anclas}). El script \cmd{indiceCurva.m} hace exactamente eso con miles de puntos.

\begin{ejemplo}[label=ej:oc-anclas]{Dos cálculos que anclan todo lo demás}
Calcule el índice del círculo unitario para el nodo $f(x)=-x$ y para la silla $f(x)=(x_1,-x_2)$.
\solucion Sobre el círculo, $x=(\cos s,\sin s)$, $s\in[0,2\pi]$. \emph{Nodo:} $f=(-\cos s,-\sin s)$, de ángulo $\theta_f=s+\pi$; cuando $s$ avanza $2\pi$, $\theta_f$ avanza $2\pi$: $I=+1$. \emph{Silla:} $f=(\cos s,-\sin s)$, de ángulo $\theta_f=-s$; cuando $s$ avanza $2\pi$, $\theta_f$ \emph{retrocede} $2\pi$: $I=-1$. En la silla, el campo gira al revés que el punto que recorre la curva (\cref{fig:oc-anclas}). Con cuentas idénticas, un foco ($f=Ax$ con $A=\begin{bmatrix}1&-1\\1&1\end{bmatrix}$) y un centro ($f=(-x_2,x_1)$) dan $+1$.
\end{ejemplo}

\video{V5_2_indice}{El vector da vueltas: nodo, silla, $z^2$}
\begin{figure}[ht]
\centering
\includegraphics[width=\linewidth]{cap05/f_indice_anclas}
\caption{Índice del círculo unitario para un nodo (arriba) y una silla (abajo). Columna 1: el campo en doce puntos numerados. Columna 2: los mismos vectores trasladados al origen. Columna 3: el ángulo acumulado, en vueltas.}
\label{fig:oc-anclas}
\end{figure}

¿Por qué esos dos números alcanzan para casi todo? Porque el índice es un \clave{invariante topológico}: si $\Gamma$ se deforma continuamente sin pasar por ningún equilibrio, el índice no cambia. La razón es corta: el índice depende continuamente de la curva (el campo es continuo y no se anula sobre ella) y solo toma valores enteros; una función continua con valores enteros es constante. En particular, el índice de un círculo pequeño alrededor de un equilibrio aislado no depende del radio: es el \clave{índice del equilibrio}, $I_f(\eq)$.

\begin{proposicion}[label=prop:oc-signo]{El índice se lee en el jacobiano}
Si $\eq$ es un equilibrio con $\det\jacat fx{\eq}\neq0$, entonces $I_f(\eq)=\sgn\det\jacat fx{\eq}$.
\end{proposicion}

\begin{proof}
Sea $A=\jacat fx{\eq}$. En un círculo de radio $\epsilon$ alrededor de $\eq$, $f(x)=A(x-\eq)+o(\epsilon)$ y $\norm{A(x-\eq)}\ge\epsilon/\norm{A^{-1}}$, así que la familia $A(x-\eq)+s\bigl(f(x)-A(x-\eq)\bigr)$, $s\in[0,1]$, no se anula sobre el círculo si $\epsilon$ es pequeño: $f$ y $A(x-\eq)$ tienen el mismo índice. Para el campo lineal, las matrices con $\det A>0$ se pueden deformar continuamente, sin pasar por $\det=0$, hasta $-I$ (el nodo, índice $+1$); las de $\det A<0$, hasta $\diag(1,-1)$ (la silla, índice $-1$).
\end{proof}

Así, nodos, focos y centros (de la linealización) tienen índice $+1$ y las sillas $-1$: es lo que calculamos en el \cref{ej:oc-anclas}, y ahora sabemos por qué. Si $\det\jacat fx{\eq}=0$, el índice puede ser cualquier entero: el campo $f=(x_1^2-x_2^2,\,2x_1x_2)$ (en complejos, $z\mapsto z^2$) tiene índice $+2$ en el origen.

%% ---------------------------------------------------------------------
\subsection{Propiedades y el teorema}
%% ---------------------------------------------------------------------

\begin{enumerate}[label=(\roman*)]
  \item \emph{Si $\Gamma$ no encierra equilibrios, $I_f(\Gamma)=0$.} Se puede encoger $\Gamma$ hasta un punto sin cruzar equilibrios; en una curva diminuta $f$ es casi constante y no da vueltas.
  \item \emph{Aditividad.} Si $\Gamma$ encierra los equilibrios aislados $x_1^*,\dots,x_k^*$, entonces $I_f(\Gamma)=\sum_iI_f(x_i^*)$. Se deforma $\Gamma$ en $k$ círculos pequeños unidos por puentes recorridos dos veces en sentidos opuestos, que se cancelan.
  \item \emph{Invertir el tiempo no cambia el índice.} Cambiar $f$ por $-f$ suma $\pi$ a todos los ángulos, y el cambio neto no se altera. Un nodo estable y uno inestable tienen el mismo índice.
  \item \emph{Una órbita cerrada tiene índice $+1$.} Sobre una órbita, $f$ es tangente a la curva, y la tangente de una curva cerrada simple da exactamente una vuelta (el \emph{Umlaufsatz} de Hopf \citep{Hop35}).
\end{enumerate}

\begin{teorema}[label=thm:oc-indice]{Lo que encierra una órbita cerrada}
Si $\gamma$ es una órbita cerrada de $\dot x=f(x)$, $x\in\R^2$, y los equilibrios que encierra son aislados, la suma de sus índices es $+1$. En consecuencia:
\begin{enumerate}
  \item toda órbita cerrada encierra al menos un equilibrio;
  \item si encierra un solo equilibrio, este no es una silla;
  \item \textbf{si todos los equilibrios encerrados tienen $\det\jacat fx{\eq}\neq0$}, y encierra $N$ nodos, focos o centros y $S$ sillas, entonces $N=S+1$ \citep[\S2.6]{Kha02}.
\end{enumerate}
\end{teorema}

\begin{proof}
Por (iv), $I_f(\gamma)=+1$; por (i) y (ii), es la suma de los índices encerrados. Los corolarios: si no encerrara nada, el índice sería $0$; un solo equilibrio de índice $-1$ daría $-1$; con equilibrios no degenerados, $N\cdot(+1)+S\cdot(-1)=1$.
\end{proof}

\begin{ideaclave}
El teorema da una condición \textbf{necesaria}: si los equilibrios dentro de una curva no suman $+1$, esa curva no puede ser una órbita. No da una condición suficiente: un nodo estable solo cumple la cuenta y no tiene órbitas cerradas alrededor.
\end{ideaclave}

\begin{ejemplo}[label=ej:oc-indicekhalil]{¿Qué puede rodear una órbita?}
Para $\dot x_1=-x_1+x_1x_2$, $\dot x_2=x_1+x_2-2x_1x_2$, diga qué combinaciones de equilibrios puede rodear una órbita cerrada.
\solucion Equilibrios: $x_1(x_2-1)=0$; con $x_1=0$ sale $x_2=0$, y con $x_2=1$ sale $x_1=1$. El jacobiano es $\begin{bmatrix}-1+x_2&x_1\\1-2x_2&1-2x_1\end{bmatrix}$. En $(0,0)$: $\begin{bmatrix}-1&0\\1&1\end{bmatrix}$, valores propios $\pm1$, $\det=-1$: silla, índice $-1$. En $(1,1)$: $\begin{bmatrix}0&1\\-1&-1\end{bmatrix}$, valores propios $-\tfrac12\pm j\tfrac{\sqrt3}2$, $\det=1$: foco estable, índice $+1$. Una órbita que rodeara a los dos tendría índice $0$; una que rodeara solo la silla, $-1$. La única posibilidad es una órbita que rodee solo el foco (\cref{fig:oc-indiceej}). El índice no dice si esa órbita existe.
\end{ejemplo}

\begin{figure}[ht]
\centering
\includegraphics[width=0.66\linewidth]{cap05/f_indice_ejemplo}
\caption{Campo del \cref{ej:oc-indicekhalil}. Una curva que rodea la silla y el foco tiene índice $0$ y no puede ser una órbita; una que rodea solo el foco tiene índice $+1$.}
\label{fig:oc-indiceej}
\end{figure}

\codigorepo{cap05_orbitas_ciclos/C5_Ej7_Indice.m}{ejemplo}{Índices del ejemplo, numéricamente}

\begin{ejemplo}[label=ej:oc-degenerado]{Cuando la linealización no ve nada, el índice sí}
Muestre que $\dot x_1=x_1^2x_2$, $\dot x_2=x_1^2-x_2^2$ no tiene \emph{ninguna} órbita cerrada.
\solucion El único equilibrio es el origen ($x_1^2x_2=0$ y $x_1^2=x_2^2$ obligan a $x_1=x_2=0$), y su jacobiano es la matriz nula: la linealización no dice nada y la \cref{prop:oc-signo} no aplica. Calculamos el índice en el círculo unitario: en $(1,0)$, $f=(0,1)$; en $\tfrac1{\sqrt2}(1,1)$, $f=\tfrac1{2\sqrt2}(1,0)$; siguiendo así, el vector gira $180^\circ$ en sentido horario, luego $360^\circ$ en sentido antihorario y otra vez $180^\circ$ en sentido horario: $[\theta_f]=-\pi+2\pi-\pi=0$. El origen tiene índice $0$. Una órbita cerrada tendría que encerrar equilibrios con índice total $+1$, y el único que hay tiene índice $0$. No hay órbitas cerradas en ninguna parte del plano.
\end{ejemplo}

\codigorepo{cap05_orbitas_ciclos/C5_Ej7_Indice.m}{indices}{Índices de nodo, foco, centro, silla, $z^2$ y el caso degenerado}

\begin{trampa}[label=tr:oc-indice]{El índice no garantiza la órbita}
\plausible{La curva $\Gamma$ encierra un foco, de índice $+1$; por lo tanto hay una órbita cerrada alrededor del foco.}
El \cref{thm:oc-indice} va en una sola dirección: órbita cerrada $\Rightarrow$ índice $+1$. El recíproco es falso: $\dot x=-x$ tiene un nodo de índice $+1$ y ninguna órbita cerrada. En el \cref{ej:oc-indicekhalil}, el índice permite una órbita alrededor del foco estable, pero no la produce.
\correcto{Use el índice para \emph{descartar} (``esta curva no puede ser una órbita''). Para \emph{probar} que una órbita existe hace falta otra herramienta: Poincaré--Bendixson (\cref{sec:oc-pb}) o Liénard.}
\end{trampa}

\begin{cuidado}[La hipótesis de $N=S+1$]
La fórmula $N=S+1$ solo vale si los equilibrios encerrados son no degenerados. El \cref{ej:oc-degenerado} lo muestra: su único equilibrio no es nodo, ni foco, ni centro, ni silla, y la fórmula no dice nada. La versión que siempre vale es la suma de índices igual a $+1$.
\end{cuidado}

\begin{profundiza}[Poincaré--Hopf: por qué el número es $+1$]
El teorema de Poincaré--Hopf dice que, para un campo en una superficie compacta (que en el borde apunta hacia afuera, si hay borde), la suma de los índices de los equilibrios es la característica de Euler $\chi$ de la superficie. El disco tiene $\chi=1$: por eso una órbita, que hace de borde de un disco invariante, encierra índice total $+1$. El círculo tiene $\chi=0$: por eso en $\dot\theta=\omega-a\sin\theta$ los equilibrios aparecen en pares estable--inestable (índices $+1$ y $-1$ en una dimensión). El toro tiene $\chi=0$: por eso dos osciladores de fase pueden derivar sin ningún equilibrio, como en la \cref{sec:oc-acoplados}.
\end{profundiza}

%% ---------------------------------------------------------------------
\section{Cómo descartar órbitas cerradas}
\label{sec:oc-descartar}
%% ---------------------------------------------------------------------

El índice descarta curvas que encierran los equilibrios equivocados. Hay tres certificados más, que descartan órbitas en regiones enteras. Los tres tienen la misma lógica: una cantidad que, a lo largo de una órbita cerrada, tendría que volver a su valor inicial, y que el sistema obliga a cambiar.

\begin{teorema}[label=thm:oc-gradiente]{Sistemas gradiente}
Si $\dot x=-\nabla V(x)$ con $V\in\Cont1$, el sistema no tiene órbitas cerradas (en cualquier dimensión).
\end{teorema}

\begin{proof}
A lo largo de una solución, $\Vdot=\nabla V\T\dot x=-\norm{\nabla V}^2\le0$, con igualdad solo en equilibrios. Si $x(t)$ fuera periódica no constante con período $T$, entonces $0=V(x(T))-V(x(0))=-\int_0^T\norm{\nabla V}^2\,\dif t<0$.
\end{proof}

\begin{ejemplo}[label=ej:oc-gradiente]{Reconocer un gradiente}
Muestre que $\dot x_1=\sin x_2$, $\dot x_2=x_1\cos x_2$ no tiene órbitas cerradas.
\solucion Con $V(x)=-x_1\sin x_2$: $-\partial V/\partial x_1=\sin x_2$ y $-\partial V/\partial x_2=x_1\cos x_2$. Es un sistema gradiente. (Una prueba rápida: un campo plano es un gradiente en $\R^2$ si y solo si $\partial f_1/\partial x_2=\partial f_2/\partial x_1$; aquí ambas valen $\cos x_2$.)
\end{ejemplo}

\begin{teorema}[label=thm:oc-lyapunov]{Una función que siempre decrece}
Si existe $V\in\Cont1$ con $\Vdot(x)=\nabla V(x)\T f(x)<0$ para todo $x$ que no sea equilibrio, el sistema no tiene órbitas cerradas.
\end{teorema}

\noindent La prueba es la misma: a lo largo de una órbita de período $T$, $V$ tendría que bajar estrictamente y volver a su valor. Note que el teorema \emph{no} pide que $V$ sea definida positiva ni radialmente no acotada.

\begin{ejemplo}[label=ej:oc-lyapunov]{Una función de Lyapunov como certificado}
Muestre que $\dot x_1=-x_1+4x_2$, $\dot x_2=-x_1-x_2^3$ no tiene órbitas cerradas, con $V=x_1^2+4x_2^2$.
\solucion $\Vdot=2x_1(-x_1+4x_2)+8x_2(-x_1-x_2^3)=-2x_1^2-8x_2^4$, que es negativa salvo en el origen (el único equilibrio). Los términos cruzados $8x_1x_2$ se cancelan: el peso $4$ en $V$ se escogió para eso.
\end{ejemplo}

\codigorepo{cap05_orbitas_ciclos/C5_Ej8_Descartar.m}{lyapunov}{$\Vdot<0$ fuera del origen, verificado en una malla}

\begin{trampa}[label=tr:oc-gas]{Para descartar órbitas no hace falta estabilidad global}
\plausible{Para descartar órbitas cerradas con una función $V$ hay que probar que el origen es globalmente asintóticamente estable: $V$ definida positiva, radialmente no acotada y $\Vdot<0$.}
Esas hipótesis dan GAS (\cref{cap:lyapunov}), y GAS implica que no hay órbitas cerradas; pero para descartar órbitas basta mucho menos. El \cref{thm:oc-lyapunov} solo usa que $V$ baja estrictamente fuera de los equilibrios: $V$ puede ser indefinida, y el sistema puede tener varios equilibrios, algunos inestables. En un sistema gradiente con dos mínimos y una silla, $V$ no es definida positiva respecto a ningún punto, y aun así descarta todas las órbitas.
\correcto{Una cantidad que decrece estrictamente a lo largo de toda trayectoria no estacionaria prohíbe las órbitas cerradas. Es la misma idea que usaremos con la entropía relativa del replicador y con la aptitud media en la \cref{sec:oc-nowak}.}
\end{trampa}

%% ---------------------------------------------------------------------
\subsection{El criterio de Bendixson--Dulac}
%% ---------------------------------------------------------------------

\begin{teorema}[label=thm:oc-bendixson]{Criterio de Bendixson--Dulac \citep{Ben01,Dul37}}
Sea $D\subseteq\R^2$ una región \emph{simplemente conexa} (sin huecos), $f\in\Cont1$ y $g:D\to\R$ una función $\Cont1$. Si
\[
  \operatorname{div}(gf)=\frac{\partial(gf_1)}{\partial x_1}+\frac{\partial(gf_2)}{\partial x_2}
\]
no es idénticamente cero y no cambia de signo en $D$, entonces $D$ no contiene órbitas cerradas. Con $g\equiv1$ es el criterio de Bendixson.
\end{teorema}

\begin{proof}
Suponga que $\gamma\subset D$ es una órbita cerrada y sea $S$ la región que encierra; $S\subset D$ porque $D$ es simplemente conexa. Sobre $\gamma$, $\dif x_1=f_1\dif t$ y $\dif x_2=f_2\dif t$, así que $g(f_1\dif x_2-f_2\dif x_1)=g(f_1f_2-f_2f_1)\dif t=0$. Por el teorema de Green,
\[
  0=\oint_\gamma g\,(f_1\,\dif x_2-f_2\,\dif x_1)=\iint_S\operatorname{div}(gf)\,\dif x_1\dif x_2 .
\]
Una función continua de un solo signo, que no es idénticamente cero, no puede tener integral cero sobre $S$.
\end{proof}

\begin{ejemplo}[label=ej:oc-duffingfriccion]{La fricción destruye las órbitas de Duffing}
Muestre que $\ddot x+\delta\dot x+x^3-x=0$ no tiene órbitas cerradas si $\delta\neq0$.
\solucion Con $f=(x_2,\ x_1-x_1^3-\delta x_2)$, $\operatorname{div}f=0-\delta=-\delta$, de signo constante en todo $\R^2$. Con $\delta=0$ la divergencia es cero y Bendixson no dice nada, en concordancia con la familia de órbitas de la \cref{fig:oc-duffing}.
\end{ejemplo}

\codigorepo{cap05_orbitas_ciclos/C5_Ej8_Descartar.m}{bendixson}{Divergencia de Duffing y de Lotka--Volterra con el factor de Dulac}

La función $g$ de Dulac es un grado de libertad que hay que adivinar. Un caso clásico: en un modelo de poblaciones $\dot v=vF(v,p)$, $\dot p=pG(v,p)$, el factor $g=1/(vp)$ cancela los prefactores y da $\operatorname{div}(gf)=\frac1p\partial F/\partial v+\frac1v\partial G/\partial p$, que muchas veces tiene signo fijo en el cuadrante positivo (\cref{sec:oc-lotka}).

\begin{trampa}[label=tr:oc-bendixson]{Bendixson solo descarta}
\plausible{$\operatorname{div}f$ cambia de signo en $D$; por lo tanto $D$ contiene una órbita cerrada.}
El criterio es una condición suficiente para que \emph{no} haya órbitas; si la divergencia cambia de signo, no dice nada. El sistema gradiente $\dot x_1=-x_1^3$, $\dot x_2=x_2^3$ (con $V=\tfrac14(x_1^4-x_2^4)$) tiene $\operatorname{div}f=3(x_2^2-x_1^2)$, que cambia de signo, y no tiene ninguna órbita cerrada (\cref{thm:oc-gradiente}). En el otro sentido, el ciclo del \cref{ej:oc-polar} vive en una región donde $\operatorname{div}f=2-4r^2$ cambia de signo, como debe ser: sobre el ciclo, la integral de la divergencia en el disco que encierra tiene que ser cero.
\correcto{Bendixson--Dulac solo se usa para descartar. Para probar que hay una órbita hace falta Poincaré--Bendixson. Y si $D$ tiene un hueco (un anillo alrededor de un equilibrio), el criterio no se puede usar: la región encerrada por la órbita no está en $D$.}
\end{trampa}

%% ---------------------------------------------------------------------
\section{Cómo probar que existe una órbita: Poincaré--Bendixson}
\label{sec:oc-pb}
%% ---------------------------------------------------------------------

Las herramientas de las dos secciones anteriores dicen dónde \emph{no} puede haber órbitas. El resultado de esta sección es el único, en dimensión dos, que permite probar que una órbita \emph{existe} sin conocerla.

\begin{teorema}[label=thm:oc-pb]{Poincaré--Bendixson \citep[Lema~2.1]{Kha02}}
Considere $\dot x=f(x)$, $x\in\R^2$, con $f$ localmente Lipschitz, y sea $M$ un conjunto cerrado y acotado tal que
\begin{enumerate}
  \item $M$ no contiene equilibrios, o contiene \emph{un solo} equilibrio, y en él el jacobiano tiene valores propios con parte real positiva (un foco o un nodo inestable);
  \item toda trayectoria que empieza en $M$ se queda en $M$ para todo $t\ge0$ ($M$ es positivamente invariante).
\end{enumerate}
Entonces $M$ contiene una órbita periódica.
\end{teorema}

\noindent\emph{Idea de la prueba.} Si $M$ no tiene equilibrios: una trayectoria que empieza en $M$ es acotada, así que su conjunto $\omega$-límite es no vacío, compacto e invariante (\cref{lem:oc-omega}) y no contiene equilibrios. En el plano, un conjunto así es una órbita cerrada; el argumento usa que una curva cerrada simple separa el plano en dos regiones (teorema de Jordan), y por eso \emph{no} se generaliza a más dimensiones. Si $M$ tiene un foco o nodo inestable, se le quita un disco pequeño alrededor de él, en cuyo borde el campo apunta hacia afuera; lo que queda es un anillo positivamente invariante sin equilibrios (\cref{fig:oc-pbkhalil}). \hfill$\square$

\begin{cuidado}[Cómo se verifica la invariancia]
Para $M=\{V(x)\le c\}$, con $V\in\Cont1$, basta que el campo no apunte hacia afuera en la frontera: $\nabla V(x)\T f(x)\le0$ en $\{V=c\}$. Para un anillo $M=\{W(x)\ge c_1,\ V(x)\le c_2\}$ hace falta además $\nabla W\T f\ge0$ en $\{W=c_1\}$. Es un producto interno con el gradiente, no un producto de vectores, y se calcula sin resolver la ecuación.
\end{cuidado}

\begin{ejemplo}[label=ej:oc-pbkhalil]{Una región que atrapa un ciclo}
Pruebe que $\dot x_1=x_1+x_2-x_1(x_1^2+x_2^2)$, $\dot x_2=-2x_1+x_2-x_2(x_1^2+x_2^2)$ tiene una órbita periódica.
\solucion \emph{Equilibrios.} $f(x)=0$ se escribe $Ax=r^2x$ con $A=\begin{bmatrix}1&1\\-2&1\end{bmatrix}$ y $r^2=x_1^2+x_2^2$: un $x\neq0$ sería un vector propio de $A$ con valor propio real $r^2$, y $A$ no tiene valores propios reales. El único equilibrio es el origen. Su jacobiano $\begin{bmatrix}1&1\\-2&1\end{bmatrix}$ tiene valores propios $1\pm j\sqrt2$: foco inestable. \emph{Región.} Con $V=x_1^2+x_2^2$,
\[
  \Vdot=2(x_1^2+x_2^2)-2x_1x_2-2(x_1^2+x_2^2)^2 .
\]
En el círculo $V=c$, $\abs{2x_1x_2}\le x_1^2+x_2^2=c$, así que $\Vdot\le2c+c-2c^2=3c-2c^2$, que es $\le0$ si $c\ge\tfrac32$ (y se alcanza: con $c=1$ hay puntos de la frontera con $\Vdot=+1$). $M=\{V\le\tfrac32\}$ es compacto, positivamente invariante y contiene un solo equilibrio, un foco inestable. Por el \cref{thm:oc-pb}, $M$ contiene una órbita periódica. La simulación muestra un ciclo límite que \emph{no} es un círculo: $r$ oscila entre $0.850$ y $1.130$, con período $4.44$.
\end{ejemplo}

\script{cap05_orbitas_ciclos/C5_Ej9_PoincareBendixson.m}
\begin{figure}[ht]
\centering
\includegraphics[width=\linewidth]{cap05/f_pb_khalil}
\caption{Izquierda: el disco $V\le3/2$ (verde) atrapa las trayectorias; al quitar un disco pequeño alrededor del foco inestable (ámbar) queda un anillo sin equilibrios que contiene el ciclo (rojo). Derecha: el máximo de $\Vdot$ en la frontera $V=c$ es $3c-2c^2$.}
\label{fig:oc-pbkhalil}
\end{figure}

\codigorepo{cap05_orbitas_ciclos/C5_Ej9_PoincareBendixson.m}{criterio}{Las dos condiciones del criterio}

\begin{trampa}[label=tr:oc-pbciclo]{Poincaré--Bendixson no da un ciclo límite}
\plausible{El anillo $M=\{c_1\le x_1^2+x_2^2\le c_2\}$ es compacto, positivamente invariante y sin equilibrios para el oscilador armónico $\dot x_1=x_2$, $\dot x_2=-x_1$; por Poincaré--Bendixson, el oscilador armónico tiene un ciclo límite.}
El teorema garantiza una \emph{órbita periódica}, no una aislada. En el oscilador armónico el anillo está lleno de órbitas cerradas, ninguna aislada. Tampoco dice cuántas hay: puede haber varias, anidadas.
\correcto{Poincaré--Bendixson da existencia de al menos una órbita periódica en $M$. Para concluir que es un ciclo límite, y que es único, hace falta más: un argumento como el de Liénard, o que la región se pueda encoger hasta una sola curva (en el \cref{ej:oc-polar}, los anillos $\{c_1\le r^2\le c_2\}$ con $c_1<1<c_2$ son invariantes para todos esos valores, y al encogerlos queda solo $r=1$).}
\end{trampa}

\begin{trampa}[label=tr:oc-pb3d]{En tres dimensiones no hay Poincaré--Bendixson}
\plausible{Una trayectoria acotada que no converge a ningún equilibrio converge a una órbita periódica.}
Es cierto en el plano y falso en $\R^3$. El sistema de Lorenz \citep{Lor63} tiene trayectorias acotadas que no convergen ni a equilibrios ni a órbitas periódicas: se mueven para siempre en un atractor extraño. La prueba de Poincaré--Bendixson usa que una curva cerrada separa el plano en un adentro y un afuera, y en $\R^3$ una curva no separa nada.
\correcto{En dimensión dos, los únicos comportamientos de largo plazo de una trayectoria acotada son: un equilibrio, una órbita cerrada o una cadena de equilibrios unidos por trayectorias. El caos necesita al menos tres dimensiones (en tiempo continuo, invariante en el tiempo). Lo usaremos en la \cref{sec:oc-nowak}.}
\end{trampa}

%% ---------------------------------------------------------------------
\section{Presas y depredadores: tres modelos, tres respuestas}
\label{sec:oc-lotka}
%% ---------------------------------------------------------------------

Esta sección aplica el arsenal completo a tres modelos de poblaciones. Se parecen mucho y se comportan de tres maneras cualitativamente distintas: un continuo de órbitas cerradas, un equilibrio que atrae todo y un ciclo límite aislado. Distinguirlos es exactamente el trabajo de las herramientas del capítulo.

\begin{historia}[Lotka, Volterra y los peces del Adriático]
Alfred J.\ Lotka (1880--1949), químico físico de formación y luego estadístico de una aseguradora en Nueva York, mostró en 1920 que un modelo de reacciones químicas podía producir oscilaciones no amortiguadas \citep{Lot20}. Vito Volterra (1860--1940), matemático italiano conocido por sus ecuaciones integrales, llegó al mismo modelo en 1926 por otra pregunta: su yerno, el biólogo marino Umberto D'Ancona, había notado que durante la Primera Guerra Mundial, cuando casi no se pescaba en el Adriático, la proporción de peces depredadores en las capturas había aumentado \citep{Vol26}. En 1931 Volterra fue uno de los pocos profesores universitarios italianos que se negaron a jurar lealtad al régimen fascista, y perdió su cátedra. La descripción moderna del modelo y sus variantes está en \citet{HoSi98}.
\end{historia}

%% ---------------------------------------------------------------------
\subsection{El modelo clásico: un continuo de órbitas}
%% ---------------------------------------------------------------------

Con $v$ la densidad de presas y $p$ la de depredadores,
\begin{equation}
  \dot v=v(a-bp),\qquad\dot p=p(-c+dv),\qquad a,b,c,d>0 .
  \label{eq:oc-lv}
\end{equation}
Sin depredadores las presas crecen exponencialmente; sin presas los depredadores se extinguen; los encuentros (proporcionales a $vp$) pasan biomasa de unas a otros. Los ejes son invariantes ($v=0$ implica $\dot v=0$), así que el cuadrante positivo también lo es. Hay dos equilibrios: el origen, que es una silla (valores propios $a$ y $-c$), y $(c/d,\ a/b)$, con valores propios $\pm j\sqrt{ac}$: la linealización no decide (\cref{tr:oc-centro}).

Lo que decide es una primera integral. Dividiendo las ecuaciones y separando variables, $(d-c/v)\dot v=(a/p-b)\dot p$, que integrada da
\begin{equation}
  K(v,p)=dv-c\ln v+bp-a\ln p=\text{constante}.
  \label{eq:oc-K}
\end{equation}
$K$ tiene un mínimo estricto en $(c/d,a/b)$ y crece hacia los bordes del cuadrante: sus curvas de nivel son cerradas y el equilibrio es un centro rodeado de órbitas (\cref{fig:oc-lotka}, izquierda). Dos consecuencias: la amplitud de la oscilación la decide la condición inicial (como en el oscilador lineal), y el \emph{promedio} sobre un período no: integrando $\frac{\dif}{\dif t}\ln v=a-bp$ sobre un período $T$, el lado izquierdo da cero y queda $\frac1T\int_0^Tp\,\dif t=\frac ab$; igualmente $\frac1T\int_0^Tv\,\dif t=\frac cd$. Esa fue la respuesta de Volterra a D'Ancona: pescar disminuye $a$ y aumenta $c$, así que \emph{menos} pesca aumenta el promedio de depredadores ($a/b$) y disminuye el de presas ($c/d$).

\codigorepo{cap05_orbitas_ciclos/C5_Ej10_LotkaVolterra.m}{clasico}{$K$ se conserva y los promedios son el equilibrio}

\begin{figure}[ht]
\centering
\includegraphics[width=\linewidth]{cap05/f_lotka}
\caption{Tres modelos presa--depredador. Izquierda: Lotka--Volterra clásico con $(a,b,c,d)=(6,3,2,1)$, curvas de nivel de $K$ alrededor de $(2,2)$. Centro: con competencia, $(a,b,c,d,e,f)=(1,1,0.5,1,0.5,0.2)$, todo converge a $(0.636,0.682)$. Derecha: Rosenzweig--MacArthur con $K=5$, un ciclo límite.}
\label{fig:oc-lotka}
\end{figure}

\begin{trampa}[label=tr:oc-lvciclo]{Las órbitas de Lotka--Volterra no son ciclos límite}
\plausible{Lotka--Volterra oscila; las oscilaciones de las poblaciones reales son ciclos límite; luego \eqref{eq:oc-lv} explica las oscilaciones reales.}
Las órbitas de \eqref{eq:oc-lv} forman una familia continua: ninguna está aislada, la amplitud depende de la condición inicial y cualquier perturbación la cambia para siempre. Peor: el modelo es \emph{estructuralmente inestable}, y el menor término adicional destruye los centros, como muestran los dos modelos siguientes.
\correcto{\eqref{eq:oc-lv} es un modelo conservativo (tiene la primera integral $K$), igual que Duffing sin fricción. Para explicar una oscilación robusta hace falta un ciclo límite, y para eso hace falta otra no linealidad.}
\end{trampa}

%% ---------------------------------------------------------------------
\subsection{Con competencia: un equilibrio globalmente atractor}
%% ---------------------------------------------------------------------

Si las presas compiten por un recurso limitado ($e>0$) y los depredadores también ($f\ge0$),
\begin{equation}
  \dot v=v(a-ev-bp),\qquad\dot p=p(-c+dv-fp) .
  \label{eq:oc-lvc}
\end{equation}
Los equilibrios en el cuadrante cerrado son $(0,0)$, $(a/e,0)$ y, si $ad-ce>0$, el de coexistencia
\[
  v^*=\frac{af+bc}{bd+ef},\qquad p^*=\frac{ad-ce}{bd+ef} .
\]
\emph{Linealización.} En $(v^*,p^*)$ la traza del jacobiano es $-(ev^*+fp^*)<0$ y el determinante es $(ad-ce)(af+bc)/(bd+ef)>0$: los dos valores propios tienen parte real negativa, un nodo o foco \emph{estable}.

\emph{Dulac descarta órbitas.} Con $g=1/(vp)$ (\cref{sec:oc-descartar}),
\[
  \operatorname{div}\Bigl(\frac{f}{vp}\Bigr)=\frac{\partial}{\partial v}\Bigl(\frac{a-ev-bp}{p}\Bigr)+\frac{\partial}{\partial p}\Bigl(\frac{-c+dv-fp}{v}\Bigr)=-\frac ep-\frac fv<0
\]
en el cuadrante abierto, que es simplemente conexo: no hay órbitas cerradas. (Con $e=f=0$ la misma cuenta da $0$, consistente con los centros del modelo clásico.)

\emph{Una función de Lyapunov.} Volterra encontró algo más fuerte. Con
\[
  W(v,p)=d\bigl(v-v^*\ln v\bigr)+b\bigl(p-p^*\ln p\bigr),\qquad
  \dot W=-de\,(v-v^*)^2-bf\,(p-p^*)^2\le0 .
\]
$W$ tiene su mínimo en $(v^*,p^*)$ y crece hacia los bordes del cuadrante. $\dot W=0$ solo donde $v=v^*$ (si $f=0$) o en el equilibrio (si $f>0$); en la recta $v\equiv v^*$, $\dot v=0$ obliga a $p=p^*$. Por el principio de invarianza (capítulo~8), toda trayectoria que empieza en el cuadrante abierto converge a $(v^*,p^*)$ (\cref{fig:oc-lotka}, centro). Con $(a,b,c,d,e,f)=(1,1,0.5,1,0.5,0.2)$, el equilibrio es $(0.636,\ 0.682)$.

\codigorepo{cap05_orbitas_ciclos/C5_Ej10_LotkaVolterra.m}{competencia}{Convergencia con competencia y la función de Volterra}

%% ---------------------------------------------------------------------
\subsection{Rosenzweig--MacArthur: un ciclo límite}
%% ---------------------------------------------------------------------

Un depredador real no come proporcionalmente a la cantidad de presas: se sacia. Con una respuesta funcional saturada (Holling tipo II) y presas con crecimiento logístico, el modelo adimensional de Rosenzweig y MacArthur \citep{RoMa63} es
\begin{equation}
  \dot v=v\Bigl(1-\frac vK\Bigr)-\frac{vp}{1+v},\qquad\dot p=p\Bigl(\frac{\beta v}{1+v}-\gamma\Bigr),
  \label{eq:oc-rm}
\end{equation}
con $K$ la capacidad de carga del ambiente y $\beta>\gamma$. El equilibrio de coexistencia tiene $v^*=\gamma/(\beta-\gamma)$, y el cálculo de la traza del jacobiano muestra que cambia de signo en
\[
  K_H=\frac{\beta+\gamma}{\beta-\gamma} .
\]
Con $\beta=2$, $\gamma=1$: $v^*=1$ y $K_H=3$. Para $K<3$ el equilibrio es un foco estable; para $K>3$ es un foco inestable. Se puede mostrar que una región acotada del cuadrante atrapa todas las trayectorias y que las trayectorias interiores no pueden acercarse a las sillas de los ejes, $(0,0)$ y $(K,0)$; con eso, Poincaré--Bendixson da una órbita periódica alrededor del foco (es el Problema~\ref{prob:oc-rm}). La simulación con $K=5$ muestra un ciclo límite en el que las presas oscilan entre $0.047$ y $4.47$ (\cref{fig:oc-lotka}, derecha). El ciclo nace en una bifurcación de Hopf (capítulo~6).

\codigorepo{cap05_orbitas_ciclos/C5_Ej10_LotkaVolterra.m}{rosenzweig}{Del equilibrio estable al ciclo al aumentar $K$}

\begin{trampa}[label=tr:oc-enriquecer]{Más recursos no estabilizan}
\plausible{Si se aumenta la capacidad de carga $K$ de las presas ---más alimento, mejor hábitat---, el ecosistema se vuelve más estable.}
En \eqref{eq:oc-rm} pasa lo contrario: para $K>K_H$ el equilibrio pierde la estabilidad y aparece un ciclo cuya amplitud crece con $K$. Con $K=5$ las presas bajan hasta $0.047$, cerca de la extinción. Es la \emph{paradoja del enriquecimiento} \citep{Ros71}.
\correcto{La estabilidad depende de la estructura del sistema, no de la abundancia. Lo que hay que mirar es la traza del jacobiano en el equilibrio (cambio de estabilidad) y, después, una región que atrape las trayectorias (existencia del ciclo).}
\end{trampa}

\begin{ideaclave}
\begin{tabular}{@{}lll@{}}
\textbf{modelo} & \textbf{herramienta que decide} & \textbf{resultado}\\
Lotka--Volterra & primera integral $K$ & continuo de órbitas (centro)\\
con competencia & Dulac con $1/(vp)$; Lyapunov $W$ & equilibrio globalmente atractor\\
Rosenzweig--MacArthur & traza $>0$ + Poincaré--Bendixson & ciclo límite
\end{tabular}
\end{ideaclave}

%% ---------------------------------------------------------------------
\section{Aplicación: ciclos límite y lenguaje}
\label{sec:oc-nowak}
%% ---------------------------------------------------------------------

Esta sección no es un ejemplo más: es un ejercicio de \emph{todas} las herramientas del capítulo sobre un modelo publicado. \citet{MiNo04} modelaron el cambio de una lengua en una población que aprende de sus padres. Hay $n$ gramáticas posibles $G_1,\dots,G_n$; $x_j\ge0$ es la fracción de la población que usa $G_j$, con $\sum_jx_j=1$. Dos hablantes con gramáticas $G_i$ y $G_j$ se entienden con un beneficio $B_{ij}$; la \emph{aptitud} de $G_i$ es su beneficio promedio y la de la población, la media:
\[
  F_i=\sum_jB_{ij}x_j=(Bx)_i,\qquad\bar F=\sum_iF_ix_i=x\T Bx .
\]
Un niño cuyo padre usa $G_i$ aprende $G_j$ con probabilidad $Q_{ij}$ (aprendizaje imperfecto si $Q_{ii}<1$; las filas de $Q$ suman $1$). La dinámica \clave{replicador--mutador} es
\begin{equation}
  \dot x_j=\sum_{i=1}^nF_i\,x_i\,Q_{ij}-\bar F\,x_j,\qquad\text{es decir,}\qquad\dot x=Q\T(F\odot x)-\bar F\,x .
  \label{eq:oc-nowak}
\end{equation}
Con $Q=I$ es la dinámica del replicador de la teoría de juegos evolutiva (\cref{sec:un-halcon}). Tomamos $n=3$, los valores del artículo, $B_{ii}=0.88$ y $B_{ij}=0.2$ ($i\neq j$), y la familia circulante
\[
  Q(c)=\begin{bmatrix}0.8-c&0.2&c\\ c&0.8-c&0.2\\ 0.2&c&0.8-c\end{bmatrix},
\]
en la que $c$ es la probabilidad de ``saltar'' a la gramática anterior en el ciclo $G_1\to G_2\to G_3\to G_1$. Las figuras 1a y 1b del artículo corresponden a $c=0.01$ y $c=0.04$.

\notamargen{\textbf{Notación.} Usamos $\bar F$ para la aptitud media, que en el artículo es $\phi$; en el curso, $\phi$ es el flujo.}
\paragraph{Paso 1: el problema es plano.} Sea $s=\sum_jx_j$. Sumando \eqref{eq:oc-nowak} sobre $j$ y usando que las filas de $Q$ suman $1$,
\[
  \dot s=\sum_iF_ix_i-\bar F\,s=\bar F\,(1-s).
\]
El plano $s=1$ es invariante y \emph{atractor} (con tasa $\bar F>0$), y en la cara $x_j=0$, $\dot x_j=\sum_iF_ix_iQ_{ij}\ge0$: el símplex $\simplex=\{x\ge0,\ s=1\}$ es compacto y positivamente invariante. Sobre él hay dos coordenadas libres, $(x_1,x_2)$: \emph{todo el arsenal plano del capítulo se aplica literalmente}.

\paragraph{Paso 2: el equilibrio interior y su índice.} Por la simetría circulante, $x^*=(\tfrac13,\tfrac13,\tfrac13)$ es un equilibrio para todo $c$, con $\bar F^*=1.28/3=0.4267$. Restringido al símplex, el jacobiano tiene valores propios con
\[
  \Re\lambda=\frac{0.092-2.94\,c}{3},\qquad\text{por ejemplo }\lambda=0.0209\pm0.1075j\ \text{con }c=0.01 .
\]
(El tercer valor propio del jacobiano completo, $-\bar F^*=-0.4267$, es el de la dirección transversal al símplex del paso~1.) Con $c=0.01$ es un foco inestable, de índice $+1$: \emph{el índice no prohíbe una órbita alrededor de él}.

\paragraph{Paso 3: sin mutación no hay ciclos.} Con $Q=I$ y $B$ simétrica, la aptitud media cumple
\[
  \dot{\bar F}=2x\T B\dot x=2\sum_ix_i(F_i-\bar F)^2\ge0,
\]
con igualdad solo en los equilibrios (es el \emph{teorema fundamental} de Fisher). $-\bar F$ es una función que decrece estrictamente fuera de los equilibrios: por el \cref{thm:oc-lyapunov}, el replicador puro no tiene órbitas cerradas. Lo que la mutación destruye no es la linealidad ---el modelo es no lineal con o sin ella---, sino esa estructura de gradiente.

\paragraph{Paso 4: el ciclo nace en una bifurcación de Hopf.} $\Re\lambda$ cambia de signo en
\[
  c^*=\frac{0.092}{2.94}=\frac{23}{735}\approx0.0313,
\]
con parte imaginaria $0.0955$ (período $\approx65.8$ al nacer). Para $c<c^*$ el foco es inestable; para $c>c^*$, estable (\cref{fig:oc-nowak}). Es la bifurcación de Hopf del capítulo~6, en un caso real.

\paragraph{Paso 5: Poincaré--Bendixson cierra el argumento.} Con $c=0.01$: el símplex es compacto y positivamente invariante (paso~1); la frontera es \emph{repelente} (la mutación alimenta las gramáticas ausentes: en $x=(10^{-6},\tfrac12,\tfrac12)$, $\dot x_1=+0.0567$); el único equilibrio interior es un foco inestable (paso~2). Quitando un disco pequeño alrededor de él y una franja delgada junto a la frontera, queda un anillo compacto, positivamente invariante y sin equilibrios. Por el \cref{thm:oc-pb} contiene una órbita periódica: el ciclo de la figura~1a del artículo está \emph{demostrado}, no solo simulado. En la simulación, $x_1$ oscila entre $0.186$ y $0.478$ con período $\approx63$.

\script{cap05_orbitas_ciclos/C5_Ej11_Nowak.m}
\begin{figure}[ht]
\centering
\includegraphics[width=\linewidth]{cap05/f_nowak}
\caption{Replicador--mutador de tres gramáticas en el símplex. Con $c=0.01$ las trayectorias de adentro y de afuera llegan a un ciclo límite; con $c=0.04$, al equilibrio $(\tfrac13,\tfrac13,\tfrac13)$. Derecha: la parte real de los valores propios cruza cero en $c^*=23/735\approx0.0313$.}
\label{fig:oc-nowak}
\end{figure}

\codigorepo{cap05_orbitas_ciclos/C5_Ej11_Nowak.m}{hopf}{Los valores propios en el símplex y la bifurcación de Hopf}

\paragraph{Paso 6: el caos exige salir del plano.} Mitchener y Nowak pasan después a cinco gramáticas, y encuentran caos. No es realismo lingüístico: con tres gramáticas el símplex es bidimensional y Poincaré--Bendixson prohíbe el caos (\cref{tr:oc-pb3d}). Con cuatro, el símplex ya es tridimensional. La hipótesis de dimensión de todos los teoremas de este capítulo es exacta.

\begin{trampa}[label=tr:oc-simulado]{Una simulación no prueba un ciclo}
\plausible{La simulación con $c=0.01$ muestra una curva cerrada a la que llegan todas las trayectorias; eso prueba que hay un ciclo límite.}
Una simulación finita no distingue un ciclo de un foco estable muy lento ni de una espiral que se cierra en un millón de vueltas. Cerca de $c^*$ la parte real es diminuta: con $c=0.03$, $\Re\lambda=0.0013$, y una espiral inestable tarda cientos de unidades de tiempo en notarse.
\correcto{La prueba es la del paso~5: región compacta positivamente invariante $+$ un solo equilibrio, foco inestable $\Rightarrow$ órbita periódica. La simulación ilustra y sugiere; Poincaré--Bendixson demuestra.}
\end{trampa}

\begin{trampa}[label=tr:oc-fisher]{No es la no linealidad la que produce el ciclo}
\plausible{El ciclo aparece porque el replicador es no lineal; la mutación solo cambia los números.}
Con $Q=I$ el replicador es igual de no lineal y no tiene ciclos: $\bar F$ crece estrictamente (paso~3). El ciclo aparece porque el aprendizaje imperfecto ($Q\neq I$) rompe la estructura que lo prohibía.
\correcto{Antes de buscar un ciclo, pregunte qué estructura lo impediría (gradiente, función que decrece, Dulac) y si el modelo la tiene. Cuando una perturbación rompe esa estructura, puede aparecer un ciclo.}
\end{trampa}

%% ---------------------------------------------------------------------
\section{Para llevar}
\label{sec:oc-cierre}
%% ---------------------------------------------------------------------

\begin{parallevar}
\begin{enumerate}
  \item \textbf{Órbita cerrada:} la trayectoria de una solución periódica no constante. No pasa por equilibrios ni se corta. En un sistema lineal vienen en familias; un \textbf{ciclo límite} es una órbita aislada, y es un fenómeno exclusivamente no lineal.
  \item \textbf{Con valores propios imaginarios puros la linealización no concluye.} Deciden una cantidad conservada (centro), una función de Lyapunov o las coordenadas polares.
  \item \textbf{Osciladores de fase:} en el círculo hay órbitas sin equilibrios. $\dot\theta=\omega-a\sin\theta$ deriva con período $2\pi/\sqrt{\omega^2-a^2}$ si $a<\omega$ y se engancha si $a>\omega$. Dos osciladores, una luciérnaga y su linterna, y cada oscilador de Kuramoto frente al campo medio obedecen esa misma ecuación.
  \item \textbf{Kuramoto:} la incoherencia pierde la estabilidad en $K_c=2/(\pi g(0))$; para Lorentz, $r=\sqrt{1-K_c/K}$. Acoplar sincroniza frecuencias, no iguala fases.
  \item \textbf{Para descartar órbitas:} índice (la suma dentro de una órbita es $+1$), sistemas gradiente, funciones que decrecen estrictamente, Bendixson--Dulac. Son condiciones necesarias.
  \item \textbf{Para probar que existen:} Poincaré--Bendixson (región compacta que atrapa, sin equilibrios o con un solo foco/nodo inestable) y Liénard (además, unicidad). Solo en el plano.
  \item \textbf{Lotka--Volterra:} un centro (primera integral); con competencia, un equilibrio globalmente atractor (Dulac y Volterra); con saturación, un ciclo límite (Hopf y Poincaré--Bendixson). Nowak: la mutación rompe la estructura de Fisher y aparece un ciclo, demostrado con Poincaré--Bendixson.
\end{enumerate}
\end{parallevar}

\begin{trampa}[title={\faBug\hspace{0.45em}Las trampas del capítulo, en una caja}]{}
\begin{itemize}
  \item $\pm j\omega$ en la linealización no es un centro (\cref{tr:oc-centro}); un ciclo estable no es una solución asintóticamente estable (\cref{tr:oc-orbital}).
  \item Arrastrada no quiere decir en fase (\cref{tr:oc-enfase}); acoplar no es igualar (\cref{tr:oc-kuramoto}).
  \item El índice no garantiza la órbita (\cref{tr:oc-indice}); para descartar órbitas no hace falta estabilidad global (\cref{tr:oc-gas}); Bendixson solo descarta (\cref{tr:oc-bendixson}).
  \item Poincaré--Bendixson no da un ciclo \emph{límite} (\cref{tr:oc-pbciclo}) y no vale en tres dimensiones (\cref{tr:oc-pb3d}).
  \item Las órbitas de Lotka--Volterra no son ciclos límite (\cref{tr:oc-lvciclo}); más recursos no estabilizan (\cref{tr:oc-enriquecer}).
  \item Una simulación no prueba un ciclo (\cref{tr:oc-simulado}); no es la no linealidad la que produce el ciclo de Nowak (\cref{tr:oc-fisher}).
\end{itemize}
\end{trampa}

\begin{profundiza}[Correcciones respecto a las versiones anteriores (Handout4a, LotkaVolterra y las diapositivas de 2026)]
(i) El criterio de Bendixson \emph{descarta} órbitas y Poincaré--Bendixson \emph{prueba} su existencia; el handout decía exactamente lo contrario. (ii) El criterio de Poincaré--Bendixson permite \emph{un solo} equilibrio en $M$ (foco o nodo inestable), no ``al menos uno''. (iii) En el ejemplo de índice, $J(1,1)=\begin{bmatrix}0&1\\-1&-1\end{bmatrix}$; la versión anterior tenía una fila equivocada. (iv) $N=S+1$ requiere equilibrios no degenerados; se agregó el ejemplo degenerado de índice $0$ y su conclusión (no hay órbitas en ninguna parte). (v) En el modelo con competencia, la traza del jacobiano en el equilibrio de coexistencia es \emph{negativa}; el texto anterior decía que los valores propios tenían parte real positiva y luego llamaba estable al equilibrio. (vi) Las coordenadas polares valen para todo $r>0$, sin excluir $x_1=0$, y en el ejemplo $\dot\theta=+1$; se corrigieron $\arctan(x_2/x_2)$ y el signo de $\dot\theta$. (vii) Las ecuaciones de estado de Rayleigh tenían un término perdido; la ecuación es de 1883. (viii) Huygens observó la sincronización en 1665 (mediados, no principios, del siglo~XVII) y la publicó en 1673; el modelo sinusoidal con pesos $a_{ij}$ generales no es el de Winfree, que propuso otro acoplamiento; Kuramoto lo planteó en 1975. (ix) En la caracterización de un ciclo límite con conjuntos $\omega$-límite se cambió ``puede ser'' por la equivalencia precisa, y se agregó la distinción entre estabilidad orbital y estabilidad de la solución periódica. (x) En el modelo de Nowak, $F_i=\sum_jB_{ij}x_j$ (el índice estaba cambiado), y la referencia es \citet{MiNo04} (2004, no 1977). (xi) Se completaron los ejemplos que no tenían solución (Khalil, ej.~2.8) y las figuras pendientes (índice, Poincaré--Bendixson, Lotka--Volterra con competencia). (xii) Notación unificada: $\phi(t;x_0)$ para el flujo, $\varphi$ para diferencias de fase, $\bar F$ para la aptitud media, $j$ para la unidad imaginaria, $\nabla V\T f$ para la derivada a lo largo del campo. Números verificados con \texttt{herramientas/cap05\_numeros.py} (59 de 59) y \cmd{verificarCap05.m} (34 de 34).
\end{profundiza}

%% ---------------------------------------------------------------------
\section{Problemas}
\label{sec:oc-problemas}
%% ---------------------------------------------------------------------

Los problemas marcados con \etiquetaM{} requieren \matlab/\simulink; los marcados con \etiquetaEstrella{} son de profundización. Cada uno tiene una trampa: búsquela antes de dar la respuesta.

\subsection*{Órbitas y ciclos}

\begin{problema}[un ciclo en polares]\label{prob:oc-polar2}
Para $\dot r=r(r-1)(2-r)$, $\dot\theta=1$: (a) halle las órbitas cerradas y su estabilidad; (b) escriba el sistema en coordenadas cartesianas; (c) ¿cuál es la región de atracción del origen?
\end{problema}

\begin{problema}[el péndulo]\label{prob:oc-pendulo}
Para $\ddot\theta+\sin\theta=0$: (a) muestre que $E=\tfrac12\dot\theta^2+1-\cos\theta$ se conserva; (b) ¿son las órbitas alrededor de $(0,0)$ ciclos límite?; (c) \etiquetaM{} calcule el período en función de la amplitud. ¿Es el péndulo un oscilador isócrono? (d) Con fricción $b\dot\theta$, ¿quedan órbitas cerradas en el plano $(\theta,\dot\theta)$? ¿Y en el cilindro $S^1\times\R$?
\end{problema}

\begin{problema}[van der Pol]\label{prob:oc-vdp}
\etiquetaM{} (a) Grafique el período del ciclo de van der Pol en función de $\mu\in[0.1,10]$ y compárelo con $2\pi$ (para $\mu$ pequeño) y con $(3-2\ln2)\mu$ (para $\mu$ grande). (b) Con \cmd{crearModelo\_VanDerPol.m}, cambie el signo de $\mu$. ¿Qué tipo de ciclo queda? Explíquelo con la inversión del tiempo.
\end{problema}

\subsection*{Osciladores}

\begin{problema}[un PLL]\label{prob:oc-pll}
Un lazo de enganche de fase se modela con $\dot\varphi=\Delta\omega-K\sin\varphi$, donde $\Delta\omega$ es la diferencia entre la frecuencia de la señal y la del oscilador local. (a) ¿Cuál es el rango de enganche? (b) Con $K=2\pi\cdot100$\,rad/s y $\Delta\omega=2\pi\cdot60$\,rad/s, ¿cuál es el error de fase de régimen? (c) Si $\Delta\omega=2\pi\cdot120$\,rad/s, ¿cada cuánto se ``resbala'' un ciclo?
\end{problema}

\begin{problema}[Kuramoto con $N=3$]\label{prob:oc-kur3}
\etiquetaM{} Con $\omega=(-1,0,1)$ en \eqref{eq:oc-kuramoto}: (a) simule para $K=0.5,1,1.5,2$ y grafique $r(t)$; (b) estime el menor $K$ con el que los tres giran a la misma frecuencia; (c) ¿es $r=1$ alguna vez? Explique con la \cref{tr:oc-kuramoto}.
\end{problema}

\subsection*{Índice, criterios y Poincaré--Bendixson}

\begin{problema}[contar índices]\label{prob:oc-indice}
Para $\dot x_1=x_2$, $\dot x_2=x_1-x_1^3-0.2x_2$ (Duffing con fricción): (a) halle los equilibrios y sus índices; (b) ¿qué combinaciones podría rodear una órbita cerrada según el índice?; (c) ¿qué dice Bendixson? (d) Concluya.
\end{problema}

\begin{problema}[Dulac]\label{prob:oc-dulac}
Muestre que $\dot x=x(2-x-y)$, $\dot y=y(4x-x^2-3)$ no tiene órbitas cerradas en el cuadrante positivo. (Pista: $g=1/(xy)$.)
\end{problema}

\begin{problema}[Poincaré--Bendixson]\label{prob:oc-pb}
Para $\dot x_1=x_2$, $\dot x_2=-x_1+x_2(1-x_1^2-2x_2^2)$: (a) muestre que el origen es un foco inestable; (b) con $V=x_1^2+x_2^2$, muestre que $\Vdot\ge0$ en $V=\tfrac12$ y $\Vdot\le0$ en $V=1$; (c) concluya que hay una órbita periódica en $\tfrac12\le V\le1$. (d) \etiquetaM{} Simúlelo. ¿Es un círculo?
\end{problema}

\begin{problema}[Rosenzweig--MacArthur]\label{prob:oc-rm}
\etiquetaEstrella{} Para \eqref{eq:oc-rm}: (a) muestre que el triángulo $\{v,p\ge0,\ \beta v+p\le C\}$ es positivamente invariante para $C$ grande; (b) calcule la traza del jacobiano en el equilibrio de coexistencia y verifique $K_H=(\beta+\gamma)/(\beta-\gamma)$; (c) argumente por qué las trayectorias del interior no pueden tener como $\omega$-límite a las sillas $(0,0)$ y $(K,0)$ cuando $K>K_H$; (d) concluya la existencia del ciclo.
\end{problema}

\subsection*{La aplicación}

\begin{problema}[Nowak]\label{prob:oc-nowak}
\etiquetaM{} Con \cmd{C5\_Ej11\_Nowak.m}: (a) verifique numéricamente que $\dot s=\bar F(1-s)$ fuera del símplex; (b) grafique la amplitud del ciclo en función de $c\in[0,c^*]$ y muestre que se anula como $\sqrt{c^*-c}$ (huella de una Hopf supercrítica); (c) con $Q=I$ grafique $\bar F(x(t))$ y verifique que nunca decrece.
\end{problema}

%% ---------------------------------------------------------------------
\section{Soluciones seleccionadas}
%% ---------------------------------------------------------------------

\paragraph{Problema~\ref{prob:oc-polar2}}
(a) Órbitas en $r=1$ y $r=2$. Entre $0$ y $1$, $\dot r<0$; entre $1$ y $2$, $\dot r>0$; arriba de $2$, $\dot r<0$. Luego $r=1$ es un ciclo inestable y $r=2$ uno estable. (c) La región de atracción del origen es el disco abierto $r<1$: el ciclo inestable es su frontera.

\paragraph{Problema~\ref{prob:oc-pendulo}}
(b) No: forman una familia continua (curvas de nivel de $E$). (c) No es isócrono: el período crece con la amplitud y tiende a infinito al acercarse a la separatriz. (d) En el plano no hay órbitas cerradas ($\operatorname{div}f=-b$). En el cilindro, Bendixson no aplica (no es simplemente conexo: una órbita podría dar la vuelta al cilindro sin encerrar nada), pero la energía sí: $\dot E=-b\dot\theta^2\le0$, y no puede ser idénticamente cero a lo largo de una solución que no sea un equilibrio (si $\dot\theta\equiv0$, $\theta$ es constante y $\sin\theta=0$). En una vuelta de una órbita cerrada $E$ bajaría estrictamente y no podría volver a su valor: el argumento del \cref{thm:oc-lyapunov} descarta toda órbita cerrada, también las que dan la vuelta. Con un par constante aplicado, en cambio, el péndulo puede girar periódicamente alrededor del cilindro.

\paragraph{Problema~\ref{prob:oc-pll}}
(a) $\abs{\Delta\omega}\le K$. (b) $\varphi^*=\arcsin(0.6)\approx0.64$\,rad. (c) $T=2\pi/\sqrt{\Delta\omega^2-K^2}=2\pi/\bigl(2\pi\sqrt{120^2-100^2}\bigr)=1/66.3$\,s, unos $15$\,ms.

\paragraph{Problema~\ref{prob:oc-indice}}
(a) $(0,0)$: silla, índice $-1$; $(\pm1,0)$: focos estables, índice $+1$. (b) Una órbita podría rodear un foco solo, o los tres equilibrios ($+1-1+1=+1$). (c) $\operatorname{div}f=-0.2$: no hay órbitas cerradas en ninguna parte. (d) El índice permitía órbitas; Bendixson las descarta todas. Las dos herramientas son necesarias, no suficientes, y se complementan.

\paragraph{Problema~\ref{prob:oc-dulac}}
Con $g=1/(xy)$: $\operatorname{div}(gf)=\partial_x\bigl((2-x-y)/y\bigr)+\partial_y\bigl((4x-x^2-3)/x\bigr)=-1/y<0$.

\paragraph{Problema~\ref{prob:oc-pb}}
(b) $\Vdot=2x_2^2(1-x_1^2-2x_2^2)$. En $V=\tfrac12$: $x_1^2+2x_2^2\le2V=1$, luego $\Vdot\ge0$. En $V=1$: $x_1^2+2x_2^2\ge V=1$, luego $\Vdot\le0$. (d) No es un círculo: el ciclo vive entre las dos circunferencias sin coincidir con ninguna.

\respuestaspausas

